% Options for packages loaded elsewhere
\PassOptionsToPackage{unicode}{hyperref}
\PassOptionsToPackage{hyphens}{url}
\documentclass[
]{article}
\usepackage{xcolor}
\usepackage{amsmath,amssymb}
\setcounter{secnumdepth}{-\maxdimen} % remove section numbering
\usepackage{iftex}
\ifPDFTeX
  \usepackage[T1]{fontenc}
  \usepackage[utf8]{inputenc}
  \usepackage{textcomp} % provide euro and other symbols
\else % if luatex or xetex
  \usepackage{unicode-math} % this also loads fontspec
  \defaultfontfeatures{Scale=MatchLowercase}
  \defaultfontfeatures[\rmfamily]{Ligatures=TeX,Scale=1}
\fi
\usepackage{lmodern}
\ifPDFTeX\else
  % xetex/luatex font selection
\fi
% Use upquote if available, for straight quotes in verbatim environments
\IfFileExists{upquote.sty}{\usepackage{upquote}}{}
\IfFileExists{microtype.sty}{% use microtype if available
  \usepackage[]{microtype}
  \UseMicrotypeSet[protrusion]{basicmath} % disable protrusion for tt fonts
}{}
\makeatletter
\@ifundefined{KOMAClassName}{% if non-KOMA class
  \IfFileExists{parskip.sty}{%
    \usepackage{parskip}
  }{% else
    \setlength{\parindent}{0pt}
    \setlength{\parskip}{6pt plus 2pt minus 1pt}}
}{% if KOMA class
  \KOMAoptions{parskip=half}}
\makeatother
\usepackage{color}
\usepackage{fancyvrb}
\newcommand{\VerbBar}{|}
\newcommand{\VERB}{\Verb[commandchars=\\\{\}]}
\DefineVerbatimEnvironment{Highlighting}{Verbatim}{commandchars=\\\{\}}
% Add ',fontsize=\small' for more characters per line
\newenvironment{Shaded}{}{}
\newcommand{\AlertTok}[1]{\textcolor[rgb]{1.00,0.00,0.00}{\textbf{#1}}}
\newcommand{\AnnotationTok}[1]{\textcolor[rgb]{0.38,0.63,0.69}{\textbf{\textit{#1}}}}
\newcommand{\AttributeTok}[1]{\textcolor[rgb]{0.49,0.56,0.16}{#1}}
\newcommand{\BaseNTok}[1]{\textcolor[rgb]{0.25,0.63,0.44}{#1}}
\newcommand{\BuiltInTok}[1]{\textcolor[rgb]{0.00,0.50,0.00}{#1}}
\newcommand{\CharTok}[1]{\textcolor[rgb]{0.25,0.44,0.63}{#1}}
\newcommand{\CommentTok}[1]{\textcolor[rgb]{0.38,0.63,0.69}{\textit{#1}}}
\newcommand{\CommentVarTok}[1]{\textcolor[rgb]{0.38,0.63,0.69}{\textbf{\textit{#1}}}}
\newcommand{\ConstantTok}[1]{\textcolor[rgb]{0.53,0.00,0.00}{#1}}
\newcommand{\ControlFlowTok}[1]{\textcolor[rgb]{0.00,0.44,0.13}{\textbf{#1}}}
\newcommand{\DataTypeTok}[1]{\textcolor[rgb]{0.56,0.13,0.00}{#1}}
\newcommand{\DecValTok}[1]{\textcolor[rgb]{0.25,0.63,0.44}{#1}}
\newcommand{\DocumentationTok}[1]{\textcolor[rgb]{0.73,0.13,0.13}{\textit{#1}}}
\newcommand{\ErrorTok}[1]{\textcolor[rgb]{1.00,0.00,0.00}{\textbf{#1}}}
\newcommand{\ExtensionTok}[1]{#1}
\newcommand{\FloatTok}[1]{\textcolor[rgb]{0.25,0.63,0.44}{#1}}
\newcommand{\FunctionTok}[1]{\textcolor[rgb]{0.02,0.16,0.49}{#1}}
\newcommand{\ImportTok}[1]{\textcolor[rgb]{0.00,0.50,0.00}{\textbf{#1}}}
\newcommand{\InformationTok}[1]{\textcolor[rgb]{0.38,0.63,0.69}{\textbf{\textit{#1}}}}
\newcommand{\KeywordTok}[1]{\textcolor[rgb]{0.00,0.44,0.13}{\textbf{#1}}}
\newcommand{\NormalTok}[1]{#1}
\newcommand{\OperatorTok}[1]{\textcolor[rgb]{0.40,0.40,0.40}{#1}}
\newcommand{\OtherTok}[1]{\textcolor[rgb]{0.00,0.44,0.13}{#1}}
\newcommand{\PreprocessorTok}[1]{\textcolor[rgb]{0.74,0.48,0.00}{#1}}
\newcommand{\RegionMarkerTok}[1]{#1}
\newcommand{\SpecialCharTok}[1]{\textcolor[rgb]{0.25,0.44,0.63}{#1}}
\newcommand{\SpecialStringTok}[1]{\textcolor[rgb]{0.73,0.40,0.53}{#1}}
\newcommand{\StringTok}[1]{\textcolor[rgb]{0.25,0.44,0.63}{#1}}
\newcommand{\VariableTok}[1]{\textcolor[rgb]{0.10,0.09,0.49}{#1}}
\newcommand{\VerbatimStringTok}[1]{\textcolor[rgb]{0.25,0.44,0.63}{#1}}
\newcommand{\WarningTok}[1]{\textcolor[rgb]{0.38,0.63,0.69}{\textbf{\textit{#1}}}}
\usepackage{longtable,booktabs,array}
\newcounter{none} % for unnumbered tables
\usepackage{calc} % for calculating minipage widths
% Correct order of tables after \paragraph or \subparagraph
\usepackage{etoolbox}
\makeatletter
\patchcmd\longtable{\par}{\if@noskipsec\mbox{}\fi\par}{}{}
\makeatother
% Allow footnotes in longtable head/foot
\IfFileExists{footnotehyper.sty}{\usepackage{footnotehyper}}{\usepackage{footnote}}
\makesavenoteenv{longtable}
\setlength{\emergencystretch}{3em} % prevent overfull lines
\providecommand{\tightlist}{%
  \setlength{\itemsep}{0pt}\setlength{\parskip}{0pt}}
\usepackage{bookmark}
\IfFileExists{xurl.sty}{\usepackage{xurl}}{} % add URL line breaks if available
\urlstyle{same}
\hypersetup{
  pdftitle={ViPII: A Decree-13-Compliant Benchmark and Compact Models for Vietnamese PII Detection and Utility-Preserving De-identification},
  hidelinks,
  pdfcreator={LaTeX via pandoc}}

\title{ViPII: A Decree-13-Compliant Benchmark and Compact Models for
Vietnamese PII Detection and Utility-Preserving De-identification}
\author{}
\date{}

\begin{document}
\maketitle

\section{Section 1: Introduction}\label{section-1-introduction}

The rapid adoption of artificial intelligence in Vietnamese public
administration, healthcare, finance, and legal services has increased
the need to process documents that contain personal data. Such documents
commonly combine structured identifiers---such as citizen identity
numbers, telephone numbers, and dates of birth---with context-dependent
information, including family circumstances, health status, and
location. A system that releases these documents to an external service
for analysis may therefore create privacy, governance, and operational
risks. At the same time, simply deleting every potentially sensitive
string can make an administrative document unusable for search,
auditing, case handling, or downstream language processing.

This paper studies \textbf{utility-preserving sanitization} for
Vietnamese personal data. The objective is not only to identify
sensitive spans, but also to remove or replace them while retaining the
non-sensitive procedural content and the linguistic coherence of the
document. We focus on the terminology of Vietnam's Personal Data
Protection Decree (Decree No.~13/2023/NĐ-CP), using \texttt{PII} and
\texttt{SPI} as dataset-level shorthand for the decree's basic and
sensitive personal-data tiers, respectively. The legal mapping defines
the scope of the benchmark; it is not, by itself, a complete
legal-compliance certification for any deployment.

\subsection{1.1. Motivation and Problem
Setting}\label{motivation-and-problem-setting}

Existing multilingual PII tools provide useful starting points, but
their entity inventories and decision rules are not designed around
Vietnamese administrative conventions. Vietnamese documents contain
identity numbers with local formatting, multi-component personal names,
hierarchical addresses, and lexical items whose interpretation depends
strongly on nearby cues. For example, \emph{Nam} may denote a person's
gender, a geographic region, or part of a name, while \emph{Kinh} may
denote an ethnicity or occur inside an unrelated phrase. Sensitive
information such as health or family circumstances is often expressed as
a long narrative clause rather than a short, regular expression match.

These characteristics create two coupled technical requirements. First,
a benchmark must represent the relevant legal fields and provide
reliable character-level boundaries. Second, a sanitization system must
optimize two objectives simultaneously: reducing residual privacy
leakage and avoiding destructive over-redaction. ViPII addresses the
first requirement through controlled synthetic generation and
deterministic character-span calibration, and evaluates the second
through a dual privacy--utility metric suite.

\subsection{1.2. Research Gaps}\label{research-gaps}

We identify four gaps that motivate the study:

\begin{enumerate}
\def\labelenumi{\arabic{enumi}.}
\tightlist
\item
  \textbf{Regulatory and taxonomy mismatch.} Widely used PII resources
  and tools are generally organized around other legal or operational
  taxonomies. They do not directly expose the 34-field PII/SPI inventory
  used by this benchmark under the Vietnamese regulatory scope.
\item
  \textbf{Vietnamese administrative and linguistic specificity.} Local
  identity formats, multi-level addresses, Vietnamese name structure,
  and context-dependent words create failure cases that are
  underrepresented in general-purpose PII benchmarks.
\item
  \textbf{Reliable span construction for synthetic text.} When a
  language model is asked to generate prose and character offsets at the
  same time, changes in Unicode text and subword tokenization can
  produce incorrect boundaries. A reproducible benchmark therefore needs
  to separate text generation from coordinate calculation.
\item
  \textbf{Privacy--utility trade-off in local deployment.} Cloud-scale
  models may be difficult to use for sensitive documents because of
  governance, latency, cost, or data-residency requirements. The
  relative value of compact models adapted to the local task has not
  been established under a common Vietnamese sanitization protocol.
\end{enumerate}

\subsection{1.3. The ViPII Approach}\label{the-vipii-approach}

ViPII is a benchmark and data-construction framework for Vietnamese
PII/SPI detection and utility-preserving sanitization. Its central
design principle is to assign different responsibilities to the
generator and the annotation pipeline. A manifest fixes the sensitive
values and the fields that should appear; a language model renders those
values in natural administrative contexts; and deterministic host-side
procedures compute character spans after generation. The pipeline
combines guarded verbatim matching, cue-context disambiguation, disjoint
name-component decomposition, and semantic anchor-tag parsing. These
mechanisms construct the ground truth offline and are not treated as the
neural architecture of the evaluated models.

The resulting corpus contains 47,881 documents and 474,874 annotated
character spans covering 34 field classes. Each instance provides
natural-language input, flat non-overlapping character spans, and
sanitization targets. The benchmark reports privacy protection and
utility retention together, with document-level \texttt{FULL} success
requiring both the removal of annotated sensitive units and the
preservation of evaluated non-sensitive operational units.

\subsection{1.4. Contributions}\label{contributions}

This paper makes the following contributions:

\begin{enumerate}
\def\labelenumi{\arabic{enumi}.}
\tightlist
\item
  \textbf{A Vietnamese legal-scope benchmark.} We introduce a 34-field
  PII/SPI benchmark aligned with the operational scope of Decree
  No.~13/2023/NĐ-CP, covering structured identifiers, context-dependent
  attributes, name components, and narrative sensitive information.
\item
  \textbf{A constraint-first synthesis and annotation pipeline.} We
  present a decoupled generation process in which manifests and
  controlled profiles specify target values, while deterministic
  post-processing computes accepted character spans and resolves
  candidate conflicts.
\item
  \textbf{A joint privacy--utility evaluation framework.} We define
  sanitization and retention metrics at span, field, and document
  levels, including the strict \texttt{FULL} measure for end-to-end
  success.
\item
  \textbf{An empirical study of compact models.} On the current held-out
  benchmark, fine-tuned Qwen3 compact models substantially improve over
  their unadapted counterparts. The 0.6B model achieves 72.46\%
  \texttt{SanRec} and 70.10\% \texttt{FULL}, while the 1.7B model
  provides higher retention scores; comparisons with general-purpose
  baselines are reported under the same evaluation protocol.
\end{enumerate}

\subsection{1.5. Scope and Organization}\label{scope-and-organization}

The study focuses on Vietnamese text generated for administrative,
legal, healthcare, and related social contexts. It evaluates the quality
of synthetic annotations and the behavior of sanitization systems under
the benchmark protocol; it does not claim that synthetic data fully
represents all real-world documents or that a model score alone
establishes legal compliance. Section 2 reviews related work. Section 3
defines the tasks, legal taxonomy, and operational representation.
Section 4 describes the constraint-first synthesis and annotation
pipeline, and Section 5 reports corpus statistics and quality
diagnostics. Section 6 presents the experimental comparison and main
results. Subsequent sections analyze component effects, robustness,
error patterns, practical deployment considerations, ethics,
limitations, and reproducibility before concluding the paper.

\begin{center}\rule{0.5\linewidth}{0.5pt}\end{center}

\section{Section 2: Related Work}\label{section-2-related-work}

Research on personal-data protection in language technologies spans
entity detection, de-identification, synthetic data generation, and
privacy-aware model deployment. These lines of work are closely related
but do not solve exactly the same problem. A detection system identifies
spans that may reveal an individual; a de-identification system
transforms those spans; and a synthetic-data pipeline determines how
training examples can be created without exposing real records. ViPII
connects these perspectives for Vietnamese administrative text by
combining a legal field taxonomy, controlled generation, deterministic
character-span construction, and joint privacy--utility evaluation.

Early practical PII systems were largely built from rules, dictionaries,
and modular recognizers. Microsoft Presidio, Google Cloud Data Loss
Prevention, Philter, and Amnesia illustrate different combinations of
regular expressions, entity recognizers, dictionaries, and
general-purpose language models for finding or masking identifiers.
These tools are valuable because they support structured credentials,
configurable recognizers, and production-oriented workflows. However,
their default inventories and decision policies are not designed around
Vietnam's distinction between basic and sensitive personal data. They
also provide limited support for Vietnamese administrative formats,
including locally structured identity numbers, multi-level addresses,
compound names, and ambiguous lexical cues. A regular expression can
detect a digit sequence, but it cannot by itself determine whether that
sequence is a citizen identifier or an administrative document code;
similarly, a dictionary match for \emph{Nam} or \emph{Kinh} is
insufficient without the surrounding clause. ViPII addresses this gap by
defining an explicit 34-field inventory and by treating cue-context
disambiguation as a first-class annotation mechanism rather than as an
optional post-processing rule.

Vietnamese named-entity recognition resources provide an important
linguistic foundation but differ in scope and task objective. Benchmarks
such as VLSP NER and PhoNER have advanced Vietnamese entity recognition
for news, web, and general-domain text, while Vietnamese pretrained
encoders such as PhoBERT have enabled strong token-level sequence
labeling. These resources typically focus on proper-name categories such
as persons, organizations, locations, and miscellaneous entities.
Personal-data sanitization requires a different granularity. A telephone
number, health-insurance identifier, date of birth, or bank account may
not be a conventional named entity, and sensitive information such as a
diagnosis or family hardship can extend across an entire narrative
clause. In addition, privacy-oriented annotation must distinguish a
person's name from a non-personal occurrence of the same word and must
preserve subject roles when several individuals are mentioned in one
document. ViPII therefore uses character spans with field-level labels
and a separate legal sensitivity tier, allowing structured identifiers,
name components, contextual attributes, and open-vocabulary sensitive
clauses to coexist in one benchmark.

Synthetic data generation has emerged as a response to the data-scarcity
problem in privacy research. Because real medical, legal, financial, and
administrative records cannot generally be redistributed, recent work
has explored profile-guided generation, weak supervision, templated
synthesis, and large-scale privacy-oriented corpora. PRIVASIS
demonstrates that auxiliary control variables---such as personal
profiles, record types, and background contexts---can be used to produce
large synthetic collections and parallel sanitization data without
relying on raw private documents. Other synthetic NER and tabular-data
efforts similarly use schemas, templates, or programmatic constraints to
control the entities that appear in generated examples. These approaches
establish the feasibility of reference-free or low-risk corpus
construction, but they leave open several issues for Vietnamese legal
and administrative text: how to encode locally valid identifiers, how to
represent the distinction between basic and sensitive data, how to
generate realistic distractors, and how to recover exact spans after a
language model has rewritten the surrounding prose. ViPII builds on the
profile-guided paradigm while adding Vietnamese administrative
constraints, multi-subject scenarios, and deterministic post-generation
calibration.

Programmatic labeling and weak supervision offer complementary
strategies for reducing manual annotation cost. Frameworks such as
Snorkel combine labeling functions, heuristic rules, and conflict
resolution to produce probabilistic labels from noisy sources. Synthetic
NER construction often follows a related principle by inserting known
entities into text or by deriving labels from generation metadata. These
methods are particularly useful when a corpus contains many structured
patterns, but their labels remain dependent on the precision and
coverage of the underlying rules. ViPII adopts a stricter separation of
responsibilities: the manifest specifies which values should be
realized, the language model supplies the surrounding prose, and the
host-side pipeline computes accepted character offsets through verbatim
matching, cue-context matching, name-component decomposition, and
anchor-tag parsing. This design does not make semantic annotation
infallible; rather, it removes one major source of error---LLM-generated
coordinates---and makes the calibration procedure auditable. The
benchmark consequently distinguishes candidate spans, non-sensitive
distractors, and finalized ground-truth spans instead of treating
manifest coverage as equivalent to model recall.

The final body of related work concerns text sanitization and the
preservation of utility after private information is removed.
Traditional redaction systems replace an entity with a fixed mask or
delete the surrounding phrase, which can protect privacy at the cost of
grammaticality and task utility. More recent work frames sanitization as
a conditional generation problem: the system receives a document and a
privacy instruction, then removes, abstracts, or replaces selected
information while retaining the rest of the record. This formulation is
relevant to document search, summarization, analytics, and downstream
language processing, where a sanitized record must remain intelligible
rather than merely empty. It also motivates evaluating both residual
leakage and over-redaction. ViPII operationalizes this trade-off through
Sanitization Recall, field-level sanitization accuracy, Retention
Recall, retention accuracy, and the strict document-level \texttt{FULL}
score. In contrast to a privacy-only evaluation, the benchmark penalizes
systems that protect privacy by deleting essential administrative
content.

Compact language models provide a practical deployment point for this
problem. Large cloud models may offer broad linguistic coverage, but
sending raw Vietnamese records to an external API can conflict with
organizational governance, latency, cost, or data-residency
requirements. Small models adapted to a specific taxonomy can instead be
deployed on local servers or workstations and can be inspected,
versioned, and updated within the responsible organization. Prior work
on efficient language models and on-device privacy processing supports
this direction, but there is limited evidence for Vietnamese PII/SPI
sanitization under a common evaluation protocol. ViPII therefore
compares fine-tuned Qwen3 compact models with unadapted compact models
and general-purpose baselines, while treating the reference teacher
pipeline as an empirical upper reference rather than as a directly
comparable deployment target.

In summary, existing systems contribute strong components---industrial
recognizers, Vietnamese NER resources, synthetic-data paradigms,
weak-supervision methods, and utility-aware sanitization---but they
leave a methodological intersection insufficiently covered. ViPII
focuses on that intersection: a legally scoped Vietnamese PII/SPI
benchmark, generated from controlled profiles and administrative
contexts, annotated through deterministic character-span calibration,
and evaluated on the joint objective of privacy protection and
information retention. The next section formalizes this task, its legal
scope, and the operational representation used throughout the benchmark.

\begin{center}\rule{0.5\linewidth}{0.5pt}\end{center}

\section{Section 3: Task Definition, Legal Taxonomy, and Operational
Framework}\label{section-3-task-definition-legal-taxonomy-and-operational-framework}

In this section, we formulate the dual tasks of span-level detection and
generative sanitization in \textbf{ViPII}, ground the taxonomy in the
statutory requirements of Vietnam's Personal Data Protection Decree
(Decree No.~13/2023/NĐ-CP), and detail the deterministic operational
mechanisms and conflict-resolution policies employed for ground-truth
construction.

\begin{center}\rule{0.5\linewidth}{0.5pt}\end{center}

\subsection{3.1. Formal Task Formulation}\label{formal-task-formulation}

\subsubsection{Input Representation and Unicode Coordinate
Space}\label{input-representation-and-unicode-coordinate-space}

Let \(\mathcal{C}\) denote the finite alphabet of Unicode code points.
To ensure strict canonical equivalence across Vietnamese combining
diacritics and varying Input Method Editor (IME) representations, all
text instances are normalized to Unicode Normalization Form C (NFC). An
input document is represented as an ordered sequence of \(N\) Unicode
code points:
\[X_{\text{raw}} = (c_1, c_2, \dots, c_N), \quad c_t \in \mathcal{C}\]

Within \(X_{\text{raw}}\), sensitive personal data instances reside as
contiguous subsequences, termed \textbf{character spans}. A ground-truth
personal data span is formally denoted by the tuple:
\[y_i = \bigl(s_i, e_i, f_i, l_i\bigr)\] where: *
\(s_i \in \{0, \dots, N-1\}\) is the 0-indexed inclusive start
code-point offset in Python string representation. *
\(e_i \in \{1, \dots, N\}\) is the 0-indexed exclusive end code-point
offset (\(s_i < e_i\)). * \(f_i \in \mathcal{F}\) designates the
fine-grained attribute category selected from the \textbf{34 fields}
(\(\mathcal{F} = \{f_1, \dots, f_{34}\}\)). *
\(l_i \in \{\text{PII}, \text{SPI}\}\) denotes the statutory
classification under Decree 13: Basic Personal Data (\(\text{PII}\)) or
Sensitive Personal Data (\(\text{SPI}\)).

\begin{quote}
\textbf{Measurement Invariant}: Character offsets in ViPII are defined
strictly over \textbf{Unicode code points after NFC normalization}
(\(0 \le s_i < e_i \le \text{len}(X)\) in Python 3). They are strictly
distinct from multi-byte UTF-8 byte offsets (where Vietnamese accented
characters occupy 2--3 bytes) and from tokenizer-dependent subword BPE
token indices.
\end{quote}

\begin{verbatim}
                                  ┌─────────────────────────────────────────────────────────────┐
                                  │                     Input: Raw Text X_raw                   │
                                  └──────────────────────────────┬──────────────────────────────┘
                                                                 │
                                 ┌───────────────────────────────┴──────────────────────────────┐
                                 ▼                                                              ▼
                  ┌──────────────────────────────┐                              ┌──────────────────────────────┐
                  │           TASK 1:            │                              │           TASK 2:            │
                  │     Span-level detection     │                              │      Utility-preserving      │
                  │        and annotation        │                              │         sanitization         │
                  └──────────────┬───────────────┘                              └──────────────┬───────────────┘
                                 │                                                             │
                                 ▼                                                             ▼
                  ┌──────────────────────────────┐                              ┌──────────────────────────────┐
                  │   Character Spans Y_hat      │                              │     Sanitized Text X_san     │
                  │ {(s, e, field, PII/SPI)}     │                              │ Tag Masking / Synthetic Swap │
                  └──────────────────────────────┘                              └──────────────┬───────────────┘
                                                                                               │
                                                                 ┌─────────────────────────────┴───────────────┐
                                                                 ▼                                             ▼
                                                  ┌──────────────────────────────┐              ┌──────────────────────────────┐
                                                  │       Privacy Defense        │              │     Utility Preservation     │
                                                  │  Residual Privacy Leakage ≈0 │              │       Over-redaction ≈ 0     │
                                                  │    (High Sanitization Rec)   │              │     (High Retention Rec)     │
                                                  └──────────────────────────────┘              └──────────────────────────────┘
\end{verbatim}

\subsubsection{Clarification of Text States in the Generation
Pipeline}\label{clarification-of-text-states-in-the-generation-pipeline}

To prevent mathematical conflation between intermediate pipeline
artifacts and downstream inference inputs, we formally distinguish four
text states: 1. \(X_{\text{tagged}}\): The raw draft emitted by the
generator LLM containing inline semantic markup
\texttt{⟦field⟧...⟦/field⟧}. 2. \(X_{\text{clean}}\): The clean natural
language text resulting from parsing and completely stripping markup
tags via a deterministic LIFO stack parser; ground-truth span offsets
\(\mathcal{Y}^*\) are calibrated on this string. 3. \(X_{\text{raw}}\):
The unannotated natural language document presented as input to
downstream models (\(X_{\text{raw}} \equiv X_{\text{clean}}\) during
benchmark inference). 4. \(X_{\text{sanitized}}\): The final
de-identified text generated by the sanitization model
\(\mathcal{M}(X_{\text{raw}})\).

\subsubsection{Canonical Schema: Flat, Non-Overlapping Spans and Name
Projection}\label{canonical-schema-flat-non-overlapping-spans-and-name-projection}

The benchmark ground truth \(\mathcal{Y}^*\) exported in
\texttt{dataset.jsonl} strictly follows a \textbf{flat, non-overlapping
schema}:
\[\mathcal{Y}^* = \{y_1^*, y_2^*, \dots, y_M^*\}, \quad \text{such that } e_i^* \le s_j^* \quad \forall i < j\]

To reconcile flat sequence evaluation with the hierarchical nature of
Vietnamese civil names (where a full name encompasses a surname, middle
name, and given name), ViPII formalizes a \textbf{Disjoint Name
Projection Operator} \(\Pi_{\text{name}}\).

Let \(y_{\text{full}} = (s, e, \text{'full\_name'}, \text{'PII'})\)
denote a composite civil name span. The operator
\(\Pi_{\text{name}}(y_{\text{full}}, X)\) projects \(y_{\text{full}}\)
onto a sequence of contiguous, non-overlapping sub-spans:
\[\Pi_{\text{name}}(y_{\text{full}}, X) = \{(s_{\text{fam}}, e_{\text{fam}}, \text{'family\_name'}, \text{'PII'}), (s_{\text{mid}}, e_{\text{mid}}, \text{'middle\_name'}, \text{'PII'}), (s_{\text{giv}}, e_{\text{giv}}, \text{'given\_name'}, \text{'PII'})\}_{\text{non-empty}}\]
satisfying the boundary ordering:
\[s = s_{\text{fam}} < e_{\text{fam}} \le s_{\text{mid}} < e_{\text{mid}} \le s_{\text{giv}} < e_{\text{giv}} = e\]

The fine-grained ground-truth representation
\(\mathcal{Y}^*_{\text{decomposed}}\) is obtained via disjoint
projection:
\[\mathcal{Y}^*_{\text{decomposed}} = \Bigl(\mathcal{Y}^* \setminus \{y \in \mathcal{Y}^* \mid y.\text{field} = \text{'full\_name'}\}\Bigr) \cup \bigcup_{y \in \mathcal{Y}^*, y.\text{field} = \text{'full\_name'}} \Pi_{\text{name}}(y, X)\]

This formulation ensures that both the composite view \(\mathcal{Y}^*\)
and the decomposed view \(\mathcal{Y}^*_{\text{decomposed}}\) are
strictly flat partitions over the code-point space, avoiding invalid
nested span collisions during sequence evaluation.

\subsubsection{Task 1: Span-Level Detection and
Annotation}\label{task-1-span-level-detection-and-annotation}

In this structured prediction mode, a model
\(g: \mathcal{X} \to 2^{\mathcal{S}}\) ingests \(X_{\text{raw}}\) and
predicts a set of character spans:
\[\hat{\mathcal{Y}} = \bigl\{\hat{y}_j = (\hat{s}_j, \hat{e}_j, \hat{f}_j, \hat{l}_j)\bigr\}_{j=1}^{\hat{M}}\]

A predicted span \(\hat{y}_j\) is evaluated as a strict exact match
against a ground-truth entity \(y_i^*\) if and only if:
\[\hat{s}_j = s_i^* \quad \land \quad \hat{e}_j = e_i^* \quad \land \quad \hat{f}_j = f_i^* \quad \land \quad \hat{l}_j = l_i^*\]

\subsubsection{Task 2: Utility-Preserving
Sanitization}\label{task-2-utility-preserving-sanitization}

In this generative de-identification mode, an end-to-end model
\(\mathcal{M}: \mathcal{X} \to \mathcal{X}\) transforms
\(X_{\text{raw}}\) directly into a \textbf{sanitized text}:
\[X_{\text{sanitized}} = \mathcal{M}(X_{\text{raw}})\]

De-identification follows two operational paradigms: 1. \textbf{Tag
Masking}: Replacing sensitive spans \(X_{\text{raw}}[s_i:e_i]\) with
standardized category mask tokens:
\[\tau(f_i) \in \bigl\{\texttt{[HỌ\_TÊN]}, \texttt{[CCCD]}, \texttt{[SỐ\_ĐIỆN\_THOẠI]}, \texttt{[TÌNH\_TRẠNG\_SỨC\_KHỎE]}, \dots\bigr\}\]
2. \textbf{Consistent Synthetic Replacement}: Substituting sensitive
spans with synthetic values from an isolated demographic database,
preserving global grammatical agreement and referential coherence across
documents.

\begin{center}\rule{0.5\linewidth}{0.5pt}\end{center}

\subsection{3.2. Conceptual Objectives: Privacy Leakage versus Utility
Retention}\label{conceptual-objectives-privacy-leakage-versus-utility-retention}

The evaluation of utility-preserving de-identification balances privacy
defense against information utility. We formalize this duality
conceptually through \textbf{residual privacy leakage} and
\textbf{over-redaction}.

Let \(\mathcal{P}(X_{\text{raw}})\) denote the set of ground-truth
sensitive personal data units (spans or attributes) contained within
document \(X_{\text{raw}}\), and let \(\mathcal{U}(X_{\text{raw}})\)
denote the set of non-sensitive operational, procedural, administrative,
and syntactic information units necessary to preserve document utility.

\subsubsection{Residual Privacy Leakage}\label{residual-privacy-leakage}

When a sanitization model processes \(X_{\text{raw}}\), any personal
data entity belonging to \(\mathcal{P}(X_{\text{raw}})\) that remains
identifiable or recoverable in \(X_{\text{sanitized}}\) constitutes an
empirical privacy failure.

Let \(\mathbb{I}_{\text{leak}}(y_i^*, X_{\text{sanitized}})\) be an
indicator function evaluated under an explicit verification protocol
\(\mathcal{V}_{\text{privacy}}(y_i^*, X_{\text{sanitized}})\), returning
1 if the sensitive content of ground-truth span \(y_i^*\) persists in
\(X_{\text{sanitized}}\) in unredacted or recoverable form, and 0 if
successfully sanitized. The empirical \textbf{Residual Privacy Leakage
Rate} (\(\mathcal{L}_{\text{privacy}}\)) across a test corpus of \(D\)
documents is expressed as:

\[\mathcal{L}_{\text{privacy}} = 1 - \text{SanRec} = \frac{\sum_{d=1}^D \sum_{i=1}^{M_d} \mathbb{I}_{\text{leak}}(y_{d,i}^*, X_{d,\text{sanitized}})}{\sum_{d=1}^D M_d}\]

where \(\text{SanRec}\) denotes \textbf{Sanitization Recall}.

\begin{quote}
\textbf{Methodological Disclaimer}: Residual leakage is measured
specifically against the predefined ground-truth sensitive units under
an explicitly specified matching and privacy-audit protocol. While
\(1 - \text{SanRec}\) quantifies the empirical leakage rate of annotated
entities, in production environments privacy verification also requires
auditing partial obfuscations, model hallucinations, and linkage
attacks; these operational evaluation protocols are formalized in
Section 6. A zero-leakage target (\(\mathcal{L}_{\text{privacy}} = 0\))
is adopted as a conservative engineering objective for privacy-sensitive
deployments; it should not be interpreted as a complete legal compliance
determination.
\end{quote}

\subsubsection{Over-Redaction and Utility
Retention}\label{over-redaction-and-utility-retention}

Conversely, if the sanitization model aggressively suppresses or deletes
non-sensitive administrative content, procedural instructions, legal
references, or grammatical connectives
(\(u_k \in \mathcal{U}(X_{\text{raw}})\)), the document's downstream
utility is compromised.

Let \(\mathbb{I}_{\text{suppress}}(u_k, X_{\text{sanitized}})\) indicate
whether a valid operational information token
\(u_k \in \mathcal{U}(X_{\text{raw}})\) has been erroneously masked or
destroyed. The \textbf{Over-Redaction Rate}
(\(\mathcal{O}_{\text{redact}}\)) is defined conceptually as:

\[\mathcal{O}_{\text{redact}} = 1 - \text{RetRec} = \frac{\sum_{d=1}^D \sum_{k=1}^{K_d} \mathbb{I}_{\text{suppress}}(u_{d,k}, X_{d,\text{sanitized}})}{\sum_{d=1}^D K_d}\]

where \(\text{RetRec}\) denotes \textbf{Retention Recall}. The
operational construction of \(\mathcal{U}(X)\) (derived from procedural
form slots and lexical references), the matching policy (accounting for
semantic paraphrasing vs.~destructive deletion), and the aggregation
schemes across tokens, attributes, and records (\(\text{RetAtt}\),
\(\text{RetA/R}\), \(\text{RetRec}\)) are concretely specified in
Section 6.

\subsubsection{\texorpdfstring{Comprehensive Success Metric:
\(\text{FULL}\)}{Comprehensive Success Metric: \textbackslash text\{FULL\}}}\label{comprehensive-success-metric-textfull}

A sanitization system achieves full end-to-end success on a document if
and only if it incurs zero residual privacy leakage and zero
over-redaction of operational slots:

\[\text{FULL} = \frac{1}{D} \sum_{d=1}^D \Biggl[ \prod_{i=1}^{M_d} \bigl(1 - \mathbb{I}_{\text{leak}}(y_{d,i}^*, X_{d,\text{sanitized}})\bigr) \times \prod_{k=1}^{K_d} \bigl(1 - \mathbb{I}_{\text{suppress}}(u_{d,k}, X_{d,\text{sanitized}})\bigr) \Biggr]\]

\begin{center}\rule{0.5\linewidth}{0.5pt}\end{center}

\subsection{3.3. Legal Scope and the Decree-13 Taxonomy (34
Fields)}\label{legal-scope-and-the-decree-13-taxonomy-34-fields}

\subsubsection{Statutory Context and
Jurisdiction}\label{statutory-context-and-jurisdiction}

The statutory taxonomy of ViPII is established under the legal framework
of the Socialist Republic of Vietnam, categorized across three legal
instruments: 1. \textbf{Enacted Statutory Law (Currently in Effect)}:
\textbf{Decree No.~13/2023/NĐ-CP on Personal Data Protection (PDPD)},
promulgated by the Government of Vietnam on April 17, 2023, and
effective July 1, 2023: * \textbf{Articles 2(3) and 3}: Formally define
and list the classes comprising Basic Personal Data (\(\text{PII}\)). *
\textbf{Articles 2(4) and 4}: Formally define and list the classes
comprising Sensitive Personal Data (\(\text{SPI}\)). * \textbf{Article
8}: Stipulates strict prohibitions against illegal processing or
unauthorized disclosure of personal data. * \textbf{Article 25}: Imposes
stringent regulatory conditions on cross-border transfers of Vietnamese
citizen data, underscoring the necessity for compact, local on-premise
de-identification models. 2. \textbf{Enacted Administrative Instrument}:
\textbf{Resolution No.~202/2025/QH15} of the National Assembly on
administrative territorial restructuring, which mandates tracking
historical and merged provincial/communal designations in civil
documentation. 3. \textbf{Prospective Regulatory Horizon}: The
\textbf{Draft Law on Personal Data Protection (PDPD 2025/2026)}, which
outlines enhanced audit obligations; this draft is cited strictly as
motivational background for forward-looking compliance, not as currently
enacted statutory law.

\subsubsection{The 34 Field Classes}\label{the-34-field-classes}

Under Decree 13, personal data is partitioned into two statutory tiers:
* \textbf{Basic Personal Data (\(\text{PII}\) - 18 Fields)}: Identifiers
that uniquely or collectively identify an individual in civil,
administrative, and commercial spheres. * \textbf{Sensitive Personal
Data (\(\text{SPI}\) - 16 Fields)}: Information intimately linked to
private rights and liberties which, if breached, directly imperils
dignity, financial standing, social security, or personal safety.

\paragraph{Table 3.1: Summary of the 34 Fields under the ViPII Decree-13
Legal
Taxonomy}\label{table-3.1-summary-of-the-34-fields-under-the-vipii-decree-13-legal-taxonomy}

\emph{(For the exhaustive statutory citations, detailed legal
definitions, and canonical Vietnamese examples for all 34 fields, see
Appendix A).}

{\def\LTcaptype{none} % do not increment counter
\begin{longtable}[]{@{}
  >{\raggedright\arraybackslash}p{(\linewidth - 6\tabcolsep) * \real{0.2353}}
  >{\centering\arraybackslash}p{(\linewidth - 6\tabcolsep) * \real{0.2941}}
  >{\raggedright\arraybackslash}p{(\linewidth - 6\tabcolsep) * \real{0.2353}}
  >{\raggedright\arraybackslash}p{(\linewidth - 6\tabcolsep) * \real{0.2353}}@{}}
\toprule\noalign{}
\begin{minipage}[b]{\linewidth}\raggedright
Statutory Tier
\end{minipage} & \begin{minipage}[b]{\linewidth}\centering
Field Count
\end{minipage} & \begin{minipage}[b]{\linewidth}\raggedright
Category Clusters
\end{minipage} & \begin{minipage}[b]{\linewidth}\raggedright
Field Identifiers (\(f \in \mathcal{F}\))
\end{minipage} \\
\midrule\noalign{}
\endhead
\bottomrule\noalign{}
\endlastfoot
\textbf{Basic Personal Data}(\textbf{PII} - Decree 13, Art. 2(3) \& 3) &
\textbf{18} & \textbf{Civil \& Identity Numbers}\emph{(8 fields)} &
\texttt{cccd} (Citizen ID 12 digits), \texttt{cmnd} (Legacy ID 9
digits), \texttt{phone}, \texttt{email}, \texttt{passport\_number},
\texttt{driver\_license}, \texttt{vehicle\_plate}, \texttt{tax\_code} \\
& & \textbf{Civil Identifiers \& Names}\emph{(5 fields)} &
\texttt{full\_name}, \texttt{family\_name}, \texttt{middle\_name},
\texttt{given\_name}, \texttt{address} \\
& & \textbf{Demographic \& Digital}\emph{(5 fields)} & \texttt{dob}
(Date of birth), \texttt{gender}, \texttt{marital\_status},
\texttt{nationality}, \texttt{digital\_account} \\
\textbf{Sensitive Personal Data}(\textbf{SPI} - Decree 13, Art. 2(4) \&
4) & \textbf{16} & \textbf{Financial \& Digital Credential}\emph{(4
fields)} & \texttt{bank\_account}, \texttt{social\_insurance\_no}
(BHXH), \texttt{health\_insurance\_no} (BHYT), \texttt{eid\_credentials}
(VNeID) \\
& & \textbf{Belief, Origin \& Tracking}\emph{(6 fields)} &
\texttt{ethnicity}, \texttt{religion}, \texttt{political\_view},
\texttt{location\_data}, \texttt{behavioral\_data},
\texttt{place\_components} \\
& & \textbf{Vulnerable Life \& Health}\emph{(6 fields)} &
\texttt{health\_status}, \texttt{criminal\_record},
\texttt{private\_life}, \texttt{family\_relations},
\texttt{sexual\_orientation}, \texttt{biometric} \\
\textbf{Total} & \textbf{34} & --- & --- \\
\end{longtable}
}

\begin{center}\rule{0.5\linewidth}{0.5pt}\end{center}

\subsection{3.4. Operational Ground-Truth Construction
Mechanisms}\label{operational-ground-truth-construction-mechanisms}

\begin{quote}
\textbf{Important Architectural Clarification}: The four mechanisms
described below are \textbf{deterministic algorithmic modules executed
offline within the pipeline} to construct ground-truth annotations
(\(\mathcal{Y}^*\)) and calibrate span coordinates with zero offset
error. They do \textbf{not} constitute the neural model architecture
evaluated in Section 6. Downstream models (such as fine-tuned Qwen-3
SLMs and frontier LLMs) receive only unannotated natural language
(\(X_{\text{raw}}\)) and must learn to identify spans or generate
sanitized text end-to-end without auxiliary symbolic cues.
\end{quote}

To eliminate coordinate drift without manual re-annotation, the pipeline
partitions the 34 fields across \textbf{four primary operational
mechanisms}:

\begin{verbatim}
┌─────────────────────────────────────────────────────────────────────────────────────────────────────────┐
│                          FOUR GROUND-TRUTH CALIBRATION MECHANISMS IN VIPII                              │
├───────────────────────────────┬─────────────────────────────────────────────────────────────────────────┤
│ 1. Verbatim Matching          │ Guarded word-boundary Regex lookaround matching invariant values.       │
│    (via: verbatim - 13 fields)│ Applied to canonical credentials, phone, CCCD, email, full names.       │
├───────────────────────────────┼─────────────────────────────────────────────────────────────────────────┤
│ 2. Cue-Context Disambiguation │ 80-character sentence-bounded lookback window for trigger cues.         │
│    (via: cue_context - 10 fld)│ Disambiguates homographs (e.g., "Kinh" ethnicity vs. "kinh tế").        │
├───────────────────────────────┼─────────────────────────────────────────────────────────────────────────┤
│ 3. Name-Component             │ Algorithmic decomposition of Vietnamese personal names into             │
│    Decomposition              │ Surname, Middle Name, and Given Name components.                        │
│    (via: name_component - 3 f)│ Handled via disjoint parsing over resolved full_name spans.             │
├───────────────────────────────┼─────────────────────────────────────────────────────────────────────────┤
│ 4. Semantic Anchor-Tag        │ Linear-time O(N) LIFO stack parser decoding inline tags                 │
│    Parsing (via: anchor_tag)  │ ⟦field⟧...⟦/field⟧ for complex, free-form, variable-length SPI clauses  │
│    (8 fields: 1 PII, 7 SPI)   │ (plus conversational PII digital_account handles).                      │
└───────────────────────────────┴─────────────────────────────────────────────────────────────────────────┘
\end{verbatim}

\begin{enumerate}
\def\labelenumi{\arabic{enumi}.}
\tightlist
\item
  \textbf{Guarded Verbatim Matching (\texttt{via:\ verbatim} - 13
  fields)}:

  \begin{itemize}
  \tightlist
  \item
    \emph{Principle}: Matches invariant alphanumeric credentials and
    addresses pre-specified in the generation manifest. Uses
    Unicode-aware zero-width lookaround assertions:
    \[R(v) = \texttt{(?<!\textbackslash{}w)} + \text{re.escape}(v) + \texttt{(?!\textbackslash{}w)}\]
  \item
    \emph{Fields (13 fields: 10 PII, 3 SPI)}: Basic PII: \texttt{cccd},
    \texttt{cmnd}, \texttt{phone}, \texttt{email},
    \texttt{passport\_number}, \texttt{driver\_license},
    \texttt{vehicle\_plate}, \texttt{tax\_code}, \texttt{address},
    \texttt{full\_name}. Sensitive SPI: \texttt{bank\_account},
    \texttt{social\_insurance\_no}, \texttt{health\_insurance\_no}.
  \end{itemize}
\item
  \textbf{Cue-Context Disambiguation (\texttt{via:\ cue\_context} - 10
  fields)}:

  \begin{itemize}
  \tightlist
  \item
    \emph{Principle}: Resolves lexical homographs (e.g., ethnic
    \emph{Kinh} vs.~\emph{kinh tế}, gender \emph{Nam} vs.~geographic
    \emph{miền Nam}) and general dates. The matcher inspects a sliding
    window backwards up to \(W = 80\) characters, bounded by sentence
    terminators (\texttt{.}, \texttt{;}, \texttt{\textbackslash{}n},
    \texttt{!}, \texttt{?}). A span is extracted only when an
    authoritative trigger cue is co-located in the clause.
  \item
    \emph{Fields (10 fields: 4 PII, 6 SPI)}: Basic PII: \texttt{dob},
    \texttt{gender}, \texttt{marital\_status}, \texttt{nationality}.
    Sensitive SPI: \texttt{ethnicity}, \texttt{religion},
    \texttt{political\_view}, \texttt{location\_data},
    \texttt{behavioral\_data}, \texttt{place\_components}.
  \end{itemize}
\item
  \textbf{Name-Component Decomposition (\texttt{via:\ name\_component} -
  3 fields)}:

  \begin{itemize}
  \tightlist
  \item
    \emph{Principle}: Splits validated \texttt{full\_name} strings into
    onomastic constituents using rule-based parsing over Vietnamese
    compound surnames (\texttt{ho\_kep}) and ethnic patronymic prefixes
    (\(Y\), \(H'\)).
  \item
    \emph{Fields (3 fields: 3 PII, 0 SPI)}: Basic PII:
    \texttt{family\_name}, \texttt{middle\_name}, \texttt{given\_name}.
  \end{itemize}
\item
  \textbf{Semantic Anchor-Tag Parsing (\texttt{via:\ anchor\_tag} - 8
  fields)}:

  \begin{itemize}
  \tightlist
  \item
    \emph{Principle}: Variable-length, narrative clauses lacking fixed
    patterns are enclosed by the generator in lightweight delimiters
    \texttt{⟦field⟧...⟦/field⟧}. A single-pass \(\mathcal{O}(N)\) LIFO
    stack parser extracts exact coordinates, records ground-truth spans,
    and strips all tags to yield \(X_{\text{clean}}\).
  \item
    \emph{Fields (8 fields: 1 PII, 7 SPI)}: Basic PII (1 field):
    \texttt{digital\_account} (informal OTT conversational handles;
    standard URLs can also be extracted verbatim). Sensitive SPI (7
    fields): \texttt{health\_status}, \texttt{criminal\_record},
    \texttt{private\_life}, \texttt{family\_relations},
    \texttt{sexual\_orientation}, \texttt{biometric},
    \texttt{eid\_credentials}.
  \end{itemize}
\end{enumerate}

\subsubsection{Table 3.2: Disjoint Cross-Mapping Matrix: Decree 13
Statutory Tiers vs.~Primary Operational
Mechanisms}\label{table-3.2-disjoint-cross-mapping-matrix-decree-13-statutory-tiers-vs.-primary-operational-mechanisms}

Every field is assigned to \textbf{exactly one primary operational
mechanism}, producing an exact mathematical partition:

{\def\LTcaptype{none} % do not increment counter
\begin{longtable}[]{@{}
  >{\raggedright\arraybackslash}p{(\linewidth - 6\tabcolsep) * \real{0.2353}}
  >{\raggedright\arraybackslash}p{(\linewidth - 6\tabcolsep) * \real{0.2353}}
  >{\raggedright\arraybackslash}p{(\linewidth - 6\tabcolsep) * \real{0.2353}}
  >{\centering\arraybackslash}p{(\linewidth - 6\tabcolsep) * \real{0.2941}}@{}}
\toprule\noalign{}
\begin{minipage}[b]{\linewidth}\raggedright
Primary Operational Mechanism
\end{minipage} & \begin{minipage}[b]{\linewidth}\raggedright
Basic Personal Data (PII) {[}18 Fields{]}
\end{minipage} & \begin{minipage}[b]{\linewidth}\raggedright
Sensitive Personal Data (SPI) {[}16 Fields{]}
\end{minipage} & \begin{minipage}[b]{\linewidth}\centering
Total Partition
\end{minipage} \\
\midrule\noalign{}
\endhead
\bottomrule\noalign{}
\endlastfoot
\textbf{1. Verbatim Matching} (\texttt{via:\ verbatim}) & \texttt{cccd},
\texttt{cmnd}, \texttt{phone}, \texttt{email},
\texttt{passport\_number}, \texttt{driver\_license},
\texttt{vehicle\_plate}, \texttt{tax\_code}, \texttt{address},
\texttt{full\_name} {[}10 fields{]} & \texttt{bank\_account},
\texttt{social\_insurance\_no}, \texttt{health\_insurance\_no} {[}3
fields{]} & \textbf{13} \\
\textbf{2. Cue-Context Disambiguation} (\texttt{via:\ cue\_context}) &
\texttt{dob}, \texttt{gender}, \texttt{marital\_status},
\texttt{nationality} {[}4 fields{]} & \texttt{ethnicity},
\texttt{religion}, \texttt{political\_view}, \texttt{location\_data},
\texttt{behavioral\_data}, \texttt{place\_components} {[}6 fields{]} &
\textbf{10} \\
\textbf{3. Name-Component Decomposition}
(\texttt{via:\ name\_component}) & \texttt{family\_name},
\texttt{middle\_name}, \texttt{given\_name} {[}3 fields{]} &
\emph{(None)} {[}0 fields{]} & \textbf{3} \\
\textbf{4. Semantic Anchor-Tag Parsing} (\texttt{via:\ anchor\_tag}) &
\texttt{digital\_account} {[}1 field{]} & \texttt{health\_status},
\texttt{criminal\_record}, \texttt{private\_life},
\texttt{family\_relations}, \texttt{sexual\_orientation},
\texttt{biometric}, \texttt{eid\_credentials} {[}7 fields{]} &
\textbf{8} \\
\textbf{Total Exact Partition} & \textbf{18 Fields} & \textbf{16 Fields}
& \textbf{34 Fields} \\
\end{longtable}
}

\begin{center}\rule{0.5\linewidth}{0.5pt}\end{center}

\subsection{3.5. Conflict Resolution and Priority Lattice
Policy}\label{conflict-resolution-and-priority-lattice-policy}

During ground-truth construction, multiple matchers may yield
overlapping candidate spans. For example, an atomic credential may
appear inside a narrative \texttt{private\_life} clause, or a surname
matcher may conflict with an encompassing full name. To enforce a strict
flat representation \(\mathcal{Y}^*\), the pipeline executes a
deterministic \textbf{5-Tier Priority Lattice} conflict resolution
policy.

\subsubsection{Formal Set Relationships}\label{formal-set-relationships}

We formally partition all extracted spans into three sets: 1.
\textbf{Candidate Spans (\(\mathcal{S}_{\text{cand}}\))}: The raw set of
all span hypotheses produced by the four offline extraction mechanisms:
\[\mathcal{S}_{\text{cand}} = \mathcal{S}_{\text{sensitive-cand}} \cup \mathcal{S}_{\text{distractor}}, \quad \mathcal{S}_{\text{sensitive-cand}} \cap \mathcal{S}_{\text{distractor}} = \emptyset\]
2. \textbf{Adversarial Distractor Spans
(\(\mathcal{S}_{\text{distractor}}\))}: Non-sensitive structural tokens
injected during generation to test discrimination---specifically
12-digit administrative document codes (\texttt{doc\_code} of type
\texttt{num12}) carrying label \texttt{O}. These spans participate
actively during overlap resolution to suppress false-positive matching
against 12-digit \texttt{cccd} credentials. 3. \textbf{Resolved Set
(\(\mathcal{S}_{\text{resolved}}\))}: The maximal conflict-free subset
produced by the priority lattice operator
\(\Omega(\mathcal{S}_{\text{cand}})\). 4. \textbf{Final Ground-Truth
Spans (\(\mathcal{Y}^*\))}: The finalized, non-overlapping set of
validated PII/SPI spans exported in \texttt{dataset.jsonl}, obtained by
filtering out distractors:
\[\mathcal{Y}^* = \{s \in \mathcal{S}_{\text{resolved}} \mid \text{label}(s) \neq \text{'O'}\}, \quad \text{where } \mathcal{S}_{\text{distractor}} \cap \mathcal{Y}^* = \emptyset\]

\subsubsection{The 5-Tier Priority
Lattice}\label{the-5-tier-priority-lattice}

Candidate spans are prioritized according to semantic precision and
statutory sensitivity:
\[\text{Tier 1 (Weight 100)} \succ \text{Tier 2 (Weight 80)} \succ \text{Tier 3 (Weight 60)} \succ \text{Tier 4 (Weight 50)} \succ \text{Tier 5 (Weight 40)}\]

\begin{itemize}
\tightlist
\item
  \textbf{Tier 1 (Invariant Atomic IDs, Weight 100)}: \texttt{cccd},
  \texttt{cmnd}, \texttt{phone}, \texttt{email},
  \texttt{passport\_number}, \texttt{driver\_license},
  \texttt{vehicle\_plate}, \texttt{tax\_code}, \texttt{dob},
  \texttt{bank\_account}, \texttt{social\_insurance\_no},
  \texttt{health\_insurance\_no}, \texttt{eid\_credentials}, and non-PII
  distractor codes (\texttt{doc\_code}). \emph{Atomic credentials take
  absolute precedence and cannot be overwritten.}
\item
  \textbf{Tier 2 (Specialized SPI Clauses, Weight 80)}:
  \texttt{health\_status}, \texttt{criminal\_record},
  \texttt{biometric}, \texttt{ethnicity}, \texttt{religion},
  \texttt{political\_view}, \texttt{sexual\_orientation}.
\item
  \textbf{Tier 3 (Core Civil PII, Weight 60)}: \texttt{full\_name},
  \texttt{address}, \texttt{family\_relations},
  \texttt{digital\_account}, \texttt{marital\_status}, \texttt{gender},
  \texttt{nationality}.
\item
  \textbf{Tier 4 (Standalone Personal Names, Weight 50)}: Standalone
  \texttt{given\_name} appearing in dialogue greetings or signature
  lines outside a \texttt{full\_name} span.
\item
  \textbf{Tier 5 (Broad Contextual SPI, Weight 40)}:
  \texttt{private\_life}, \texttt{location\_data},
  \texttt{behavioral\_data}.
\end{itemize}

\subsubsection{Algorithmic Tie-Breaking and Formal Interval
Subtraction}\label{algorithmic-tie-breaking-and-formal-interval-subtraction}

When two candidate spans \(s_a, s_b \in \mathcal{S}_{\text{cand}}\)
collide
(\(s_a.\text{start} < s_b.\text{end} \land s_b.\text{start} < s_a.\text{end}\)),
collisions are resolved via a deterministic hierarchy: 1.
\textbf{Primary: Priority Weight}: Higher tier weight wins
(\(\text{weight}(s_a) > \text{weight}(s_b)\)). 2. \textbf{Secondary:
Span Length}: If weights are equal, the longer span is retained
(Longest-Match: \(e_a - s_a > e_b - s_b\)). 3. \textbf{Tertiary: Start
Offset}: If weights and lengths are identical, the earlier span wins
(Left-to-Right: \(s_a.\text{start} < s_b.\text{start}\)).

\begin{quote}
\textbf{Formal Interval Subtraction for Broad Narrative Contexts (Tier
5)}: When a Tier 5 broad narrative span
\(I_{\text{broad}} = [s_{\text{broad}}, e_{\text{broad}})\) intersects
with a set of \(K\) already accepted higher-priority intervals
\(\mathcal{I}_{\text{higher}} = \{[s_k, e_k)\}_{k=1}^K\), rather than
dropping the broad context entirely, the pipeline computes the set
difference:
\[\text{Res}(I_{\text{broad}}) = I_{\text{broad}} \setminus \bigcup_{k=1}^K [s_k, e_k) = \bigcup_{j=1}^J [s'_j, e'_j)\]
where each \([s'_j, e'_j)\) is a maximal contiguous, non-empty
sub-interval. A residual sub-interval is retained as a valid
ground-truth span if and only if \(e'_j - s'_j \ge \theta_{\text{min}}\)
(where \(\theta_{\text{min}} = 5\) code points), ensuring surrounding
domestic and narrative hardship is preserved without swallowing atomic
credentials.
\end{quote}

\begin{verbatim}
Algorithm 1: Deterministic 5-Tier Priority Overlap Resolution (Ω)
──────────────────────────────────────────────────────────────────────────────────
Input  : Raw candidate spans S_cand, Priority weight function w(·), 
         Minimum residual context threshold θ_min = 5
Output : Conflict-free ground-truth span set Y*

1:  S_sorted ← Sort S_cand by primary key w(s) descending, 
                     secondary key (e - s) descending, tertiary key s ascending
2:  S_accepted ← ∅
3:  for each span s in S_sorted do
4:      I_collide ← {a ∈ S_accepted | s.start < a.end ∧ a.start < s.end}
5:      if I_collide = ∅ then
6:          S_accepted ← S_accepted ∪ {s}
7:      else if w(s) = 40 then  // Tier 5: Broad contextual narrative span
8:          Res(s) ← [s.start, s.end) \ ⋃_{a ∈ I_collide} [a.start, a.end)
9:          for each contiguous interval [s'_j, e'_j) in Res(s) do
10:             if (e'_j - s'_j) ≥ θ_min then
11:                 s_trim ← InstantiateSpan(s'_j, e'_j, s.field, s.label, s.via)
12:                 S_accepted ← S_accepted ∪ {s_trim}
13:             end if
14:         end for
15:     end if
16: end for
17: Y* ← {s ∈ S_accepted | s.label ≠ "O"}  // Filter out adversarial non-PII distractors
18: return Sort Y* by start offset ascending
──────────────────────────────────────────────────────────────────────────────────
\end{verbatim}

By formalizing these definitions, set relations, and deterministic
policies, Section 3 provides an unambiguous, mathematically consistent
foundation for the synthesis pipeline (Section 4), benchmark statistics
(Section 5), and experimental evaluations (Section 6).

\begin{center}\rule{0.5\linewidth}{0.5pt}\end{center}

\section{Section 4: Constraint-First Data Synthesis and Annotation
Pipeline}\label{section-4-constraint-first-data-synthesis-and-annotation-pipeline}

In this section, we present the end-to-end architecture of the
\textbf{ViPII} data generation and annotation pipeline. We explain the
foundational design invariants---\emph{Controlled Synthesis by
Construction} and \emph{Deterministic Server-Side Calibration}---and
systematically trace each stage: demographic profile banking, procedural
manifest construction, adversarial distractor synthesis, two-stage
contextual prompting, deterministic markup parsing, character-span
calibration, priority lattice conflict resolution, and the resulting
parallel de-identification dataset format.

\begin{center}\rule{0.5\linewidth}{0.5pt}\end{center}

\subsection{4.1. Pipeline Overview and Core Architectural
Invariants}\label{pipeline-overview-and-core-architectural-invariants}

Existing privacy datasets frequently suffer from two methodological
failure modes: (1) relying on heuristic regular expressions and
dictionary lookups over uncurated web scrapes, which yields severe label
noise and coordinate shifts; or (2) prompting autoregressive language
models to directly predict numeric character offsets, which inevitably
triggers the \emph{Cumulative Index Offset Drift} catastrophe due to the
non-isomorphism between subword token space and multi-byte Unicode
code-point space.

To establish an authoritative benchmark aligned with Decree
No.~13/2023/NĐ-CP without exposing genuine citizen data, ViPII adopts a
\textbf{constraint-first, decoupled synthesis paradigm}. The pipeline
decouples natural language prose generation from spatial coordinate
calculation via two governing technical invariants:

\begin{enumerate}
\def\labelenumi{\arabic{enumi}.}
\tightlist
\item
  \textbf{Controlled Synthesis by Construction (Constraint-First
  Principle)}: All sensitive personal identifiers (\(\text{PII}\)) and
  sensitive attributes (\(\text{SPI}\)) are generated and locked
  deterministically at the Python runtime layer \emph{prior} to prompt
  assembly. The generative language model acts strictly as a contextual
  narrator, weaving natural administrative prose, formal petitions, and
  dialogue around fixed entity values. The model is explicitly barred
  from hallucinating new personal identity slots.
\item
  \textbf{Deterministic Server-Side Span Calibration}: Character offsets
  (\(s_i, e_i\)) and entity labels are computed entirely on the host
  server using deterministic algorithmic operators (Unicode lookaround
  matching, sentence-bounded trigger windowing, and a single-pass LIFO
  stack parser). The language model is never tasked with predicting
  numeric coordinates, entirely eliminating coordinate hallucination.
\end{enumerate}

Figure 4.1 illustrates the complete 10-node operational flow structured
across five functional columns.

\begin{Shaded}
\begin{Highlighting}[]
\NormalTok{flowchart TD}
\NormalTok{    subgraph C1["Column 1: Data Assets Initialization"]}
\NormalTok{        N1["1. Synthetic Profile Bank\textless{}br/\textgreater{}(profile\_bank.jsonl {-} 70,000 profiles)"]}
\NormalTok{        N2["2. Situational Scenario Catalog\textless{}br/\textgreater{}(scenario\_catalog.json {-} 19 scenarios)"]}
\NormalTok{        N3["3. Administrative Form Registry\textless{}br/\textgreater{}(form.json {-} 9,020 public service forms)"]}
\NormalTok{    end}

\NormalTok{    subgraph C2["Column 2: Realization \& Adversarial Preparation"]}
\NormalTok{        N4["4. Procedural Manifest Builder\textless{}br/\textgreater{}(core/manifest.py: build\_manifest)"]}
\NormalTok{        N5["5. Adversarial Distractor Injection\textless{}br/\textgreater{}(core/manifest.py: realize\_supplemental)"]}
\NormalTok{    end}

\NormalTok{    subgraph C3["Column 3: Two{-}Stage Prompt Synthesis"]}
\NormalTok{        N6["6. Contextual Prompt Builder\textless{}br/\textgreater{}(track\_a/prompts.py \& track\_b/prompts.py)"]}
\NormalTok{    end}

\NormalTok{    subgraph C4["Column 4: Generative Execution \& Verification"]}
\NormalTok{        N7["7. LLM Surface Generator\textless{}br/\textgreater{}(core/api.py: Gemini / DeepSeek)"]}
\NormalTok{        N8\{"8. Syntax \& Coverage Validator\textless{}br/\textgreater{}(core/annotation.py: coverage)"\}}
\NormalTok{        N8\_RETRY["Regeneration \& Revision Loop\textless{}br/\textgreater{}(Max 3 iterations, T=0.20)"]}
\NormalTok{    end}

\NormalTok{    subgraph C5["Column 5: Deterministic Calibration \& Output"]}
\NormalTok{        N9["9. Deterministic Span Calibrator\textless{}br/\textgreater{}(Stack Parser + Lookarounds + Lattice)"]}
\NormalTok{        N10["10. Parallel Sanitization Corpus\textless{}br/\textgreater{}(dataset.jsonl: X\_raw \& X\_sanitized)"]}
\NormalTok{    end}

\NormalTok{    N1 {-}{-}\textgreater{}|Primary \& secondary profiles| N4}
\NormalTok{    N1 {-}{-}\textgreater{}|Kinship profile borrowing| N5}
\NormalTok{    N2 {-}{-}\textgreater{}|cooccur\_fields \& spi\_targets| N4}
\NormalTok{    N2 {-}{-}\textgreater{}|supplemental\_attributes| N5}
\NormalTok{    N2 {-}.{-}\textgreater{}|Cross{-}domain keyword mapping| N3}
\NormalTok{    N3 {-}.{-}\textgreater{}|Reference metadata attachment| N10}

\NormalTok{    N4 {-}{-}\textgreater{}|Locked atomic PII \& SPI targets| N6}
\NormalTok{    N5 {-}{-}\textgreater{}|12{-}digit doc\_code \& org distractors| N6}
\NormalTok{    N6 {-}{-}\textgreater{}|Two{-}stage prompt: Outline {-}\textgreater{} Draft| N7}
\NormalTok{    N7 {-}{-}\textgreater{}|Tagged text draft X\_tagged| N8}

\NormalTok{    N8 {-}{-} "Tag syntax error / Incomplete coverage" {-}{-}\textgreater{} N8\_RETRY {-}{-}\textgreater{} N7}
\NormalTok{    N8 {-}{-} "Coverage OK \& Markup Valid" {-}{-}\textgreater{} N9}
\NormalTok{    N9 {-}{-}\textgreater{}|Clean text X\_clean \& Flat spans Y*| N10}
\end{Highlighting}
\end{Shaded}

\emph{Figure 4.1: The authoritative 10-node operational flow of the
ViPII constraint-first data synthesis and deterministic annotation
pipeline.}

\begin{center}\rule{0.5\linewidth}{0.5pt}\end{center}

\subsection{4.2. Synthetic Profile Bank and Metadata
Sources}\label{synthetic-profile-bank-and-metadata-sources}

The generation pipeline is grounded in three foundational data
repositories (\texttt{data/}), engineered to reflect the demographic and
administrative reality of Vietnam without incorporating any real-world
citizen records.

\subsubsection{\texorpdfstring{The 70,000 Artificial Profile Bank
(\texttt{profile\_bank.jsonl})}{The 70,000 Artificial Profile Bank (profile\_bank.jsonl)}}\label{the-70000-artificial-profile-bank-profile_bank.jsonl}

The demographic core comprises exactly \textbf{70,000 fully realized,
unique synthetic citizen profiles}, generated via the standalone
demographic engine \texttt{build\_profiles.py} parameterized by the
empirical census distributions in \texttt{profile\_core.json}. Each
record contains 35 personal attributes structured into 10 top-level
fields: * \texttt{profile\_id}: Unique identifier formatted as
\texttt{pf\_xxxxxx} (spanning \texttt{pf\_000000} to
\texttt{pf\_069999}). * \texttt{fields}: Complete demographic slots
covering all 18 Basic PII fields and 16 Sensitive SPI fields. *
\texttt{consistency}: Internal logic invariants ensuring strict
cross-attribute compatibility: * The birth year declared in \texttt{dob}
matches the two-digit year code embedded in the Citizen Identity Card
(\texttt{cccd}). * The provincial birth code in \texttt{cccd} matches
the primary administrative jurisdiction of the declared domicile
\texttt{address}. * Age-conditional marital distributions ensure
realistic civil statuses (e.g., divorce or widowhood are conditional on
legal adult ages). * \texttt{quasi\_key} \& \texttt{k\_estimate}:
Quasi-identifier tuples (birth year, civil gender, communal
administrative unit) with estimated local \(k\)-anonymity values
(\(k \ge 5\)) ensuring statistical uniqueness without reproducing real
individual profiles.

\paragraph{Demographic and Onomastic
Realism}\label{demographic-and-onomastic-realism}

To ensure sociolinguistic representativeness across Vietnam's
multi-ethnic population, \texttt{profile\_core.json} models onomastic
distributions spanning \textbf{all 54 officially recognized ethnic
groups}: * \textbf{Surname Distributions}: Weighted across 69 standard
single surnames (reflecting national census proportions: \emph{Nguyễn}
\textasciitilde38.4\%, \emph{Trần} \textasciitilde12.1\%, \emph{Lê}
\textasciitilde9.5\%, \emph{Phạm} \textasciitilde7.1\%) and 22
traditional compound surnames (\texttt{ho\_kep}, e.g., \emph{Nguyễn
Đình}, \emph{Trần Khắc}). * \textbf{Ethnic Minority Onomastics}:
Specific matrilineal clan prefixes and patronymic naming systems, such
as Mon-Khmer and Austronesian clan names (\emph{Rơ Châm}, \emph{Siu} for
Gia Rai; \emph{Danh}, \emph{Sơn} for Khmer; ethnic gender affixes \(Y\)
for males and \(H'\) for females among Ê Đê communities). *
\textbf{Territorial Restructuring Alignment}: Mapped across all 63
provincial codes designated by the Ministry of Public Security (MPS),
explicitly reflecting historical mergers and consolidated units under
National Assembly \textbf{Resolution No.~202/2025/QH15}.

\subsubsection{\texorpdfstring{The Situational Scenario Catalog
(\texttt{scenario\_catalog.json})}{The Situational Scenario Catalog (scenario\_catalog.json)}}\label{the-situational-scenario-catalog-scenario_catalog.json}

Public administration records rarely reveal sensitive data without
specific administrative motivation. The catalog defines \textbf{19
canonical situational scenarios} spanning civil petitions,
administrative appeals, social assistance applications, judicial
records, labor disputes, and clinical registrations. Each scenario
specifies: * \texttt{cooccur\_fields}: Canonical civil identification
fields routinely required for procedural dossiers (e.g.,
\texttt{full\_name}, \texttt{cccd}, \texttt{phone}, \texttt{address}). *
\texttt{spi\_targets}: Specific sensitive attributes naturally elicited
by the procedure (e.g., \texttt{health\_status} in medical assistance
petitions; \texttt{criminal\_record} in judicial relief declarations;
\texttt{private\_life} in poverty certification). * \texttt{register}:
Required stylistic formality, ranging from standard bureaucratic forms
to urgent narrative petitions.

\subsubsection{\texorpdfstring{Administrative Form Reference Registry
(\texttt{form.json} \&
\texttt{form\_domains.json})}{Administrative Form Reference Registry (form.json \& form\_domains.json)}}\label{administrative-form-reference-registry-form.json-form_domains.json}

To reflect genuine administrative procedure without biasing generation,
we curating \textbf{9,020 real-world public service form metadata
profiles} across ministerial domains (Justice, Public Security, Health,
Transport). Crucially, \textbf{raw administrative templates are not
pasted into generator prompts}, as verbatim templates degrade narrative
fluency and bloat context windows. Instead, form identifiers, official
procedural headings, and administrative competence tiers are injected as
\textbf{reference metadata attached to the output record}, grounding the
synthetic document within authentic legal administrative workflows.

\begin{center}\rule{0.5\linewidth}{0.5pt}\end{center}

\subsection{4.3. Procedural Manifest
Construction}\label{procedural-manifest-construction}

Before invoking the language model, the pipeline executes
\texttt{build\_manifest()} (\texttt{core/manifest.py}),
deterministically extracting and formatting the entity values that must
appear in the final text.

\begin{verbatim}
       Synthetic Profile           Situational Scenario
      (pf_012845: Citizen)        (Medical Assistance)
               │                            │
               └─────────────┬──────────────┘
                             ▼
              core/manifest.py: build_manifest()
                             │
            ┌────────────────┴────────────────┐
            ▼                                 ▼
   Primary Manifest Items             Adversarial Realization
   - cccd: "001095012345"             - doc_code: "104928374619" (num12, label: O)
   - phone: "0983123456"              - org: "UBND Phường Mai Dịch"
   - dob: "12/04/1995"                - kin_name: "Nguyễn Văn Bình" (subject: father)
   - spi: health_status [SENTINEL]
\end{verbatim}

The manifest builder performs three essential functions: 1.
\textbf{Field Selection and Sentinel Assignment}: Intersects profile
fields with scenario requirements (\texttt{cooccur\_fields} \(\cup\)
\texttt{spi\_targets}). Structured civil fields receive their concrete
string values from the profile. Free-form narrative attributes (such as
\texttt{health\_status} or \texttt{private\_life}) are assigned the
runtime token \texttt{{[}{[}LLM\_GENERATE{]}{]}}, instructing the
generator to synthesize an appropriate narrative clause within the
scenario context while wrapping it in corresponding markup tags. 2.
\textbf{Surface Variant Expansion}: Generates canonical surface variants
for structured credentials (e.g., rendering \texttt{cccd} with standard
spacing \texttt{001\ 095\ 012345} or contiguous \texttt{001095012345};
rendering phone numbers with domestic \texttt{09x} or international
\texttt{+84} prefixes) so that downstream matching algorithms account
for natural syntactic variations. 3. \textbf{Onomastic Decomposition
Preparation}: Decomposes the primary citizen name into sub-tokens
(\texttt{family\_name}, \texttt{middle\_name}, \texttt{given\_name}) via
\texttt{split\_vietnamese\_name()}, registering derived items in the
manifest to facilitate post-hoc span resolution.

\begin{center}\rule{0.5\linewidth}{0.5pt}\end{center}

\subsection{4.4. Adversarial Distractors and Multi-Subject
Realization}\label{adversarial-distractors-and-multi-subject-realization}

A common vulnerability in entity extraction and de-identification
systems is the reliance on simplistic surface regularities (e.g.,
treating any 12-digit numeric sequence as a Citizen ID). To enforce
robust contextual discrimination, \texttt{realize\_supplemental()}
(\texttt{core/manifest.py}) synthesizes three classes of adversarial
distractors:

\subsubsection{\texorpdfstring{1. 12-Digit Non-PII Administrative
Document Codes (\texttt{doc\_code} of type
\texttt{num12})}{1. 12-Digit Non-PII Administrative Document Codes (doc\_code of type num12)}}\label{digit-non-pii-administrative-document-codes-doc_code-of-type-num12}

The generator injects procedural dossier identifiers, dispatch numbers,
and receipt barcodes structured as 12-digit numeric strings:
\[\text{doc\_code} = \sum_{i=0}^{11} d_i \cdot 10^{11-i}, \quad d_i \in \{0, \dots, 9\}\]
In prompt instructions, these sequences are explicitly labeled as
procedural document identifiers
(\texttt{Mã\ hồ\ sơ\ tiếp\ nhận:\ «104928374619»}), assigned
ground-truth label \textbf{\texttt{O}} (non-PII). Models must examine
the surrounding lexical context (distinguishing \emph{``Mã hồ sơ
số\ldots{}''} from \emph{``Số định danh cá nhân\ldots{}''}) rather than
firing purely on digit count.

\subsubsection{\texorpdfstring{2. Multi-Subject Kinship Entities
(\texttt{subject:\ other})}{2. Multi-Subject Kinship Entities (subject: other)}}\label{multi-subject-kinship-entities-subject-other}

Administrative declarations frequently reference family members, legal
guardians, or guarantors. The pipeline dynamically samples secondary
profiles from the bank to inject kinship entities (e.g., a father's
name, spouse's phone number, or dependent child's birth date). These
entities are registered in the manifest with an explicit subject tag
(\texttt{subject:\ "father"}, \texttt{subject:\ "spouse"}), evaluating
whether models de-identify secondary individuals while preserving
contextual role relations.

\subsubsection{3. Affirmations of Clean Judicial
Records}\label{affirmations-of-clean-judicial-records}

Under Article 2.4(g) of Decree 13, criminal records constitute Sensitive
SPI. However, standard public administration routinely requires citizens
to state that they \emph{possess no criminal record} (e.g., \emph{``Tôi
cam đoan không có tiền án tiền sự''}). Labeling clean affirmations as
sensitive criminal records causes catastrophic over-redaction. The
manifest builder explicitly suppresses \texttt{criminal\_record}
annotations when the narrative context dictates a clean declaration,
injecting a synthetic Judicial Record Certificate serial number instead.

\begin{center}\rule{0.5\linewidth}{0.5pt}\end{center}

\subsection{4.5. Two-Stage Contextual Prompting and
Revision}\label{two-stage-contextual-prompting-and-revision}

To generate natural administrative prose while preserving 100\% manifest
adherence, generation follows a \textbf{Two-Stage Prompting Strategy}:

\begin{verbatim}
                  ┌──────────────────────────────────────────────┐
                  │            STAGE 1: OUTLINE PROMPT           │
                  │ - Citizen identity & Scenario context        │
                  │ - Procedural motivation & required structure │
                  └──────────────────────┬───────────────────────┘
                                         │
                                         ▼
                  ┌──────────────────────────────────────────────┐
                  │            STRUCTURAL OUTLINE DRAFT          │
                  │ Heading -> Administrative Body -> Disclosures│
                  └──────────────────────┬───────────────────────┘
                                         │
                                         ▼
                  ┌──────────────────────────────────────────────┐
                  │            STAGE 2: DRAFT PROMPT             │
                  │ - Full manifest entity injection             │
                  │ - Mandatory anchor tags: ⟦field⟧...⟦/field⟧  │
                  │ - Distractor integration (doc_code, org)     │
                  └──────────────────────┬───────────────────────┘
                                         │
                                         ▼
                  ┌──────────────────────────────────────────────┐
                  │            INTERMEDIATE DRAFT X_tagged       │
                  └──────────────────────┬───────────────────────┘
                                         │
                         ┌───────────────┴───────────────┐
                         ▼                               ▼
                 [Coverage OK & Valid]           [Missing / Invalid]
                         │                               │
                         │                               ▼
                         │               ┌──────────────────────────────┐
                         │               │  STAGE 2b: REVISION LOOP     │
                         │               │  Low Temperature (T = 0.20)  │
                         │               │  Fix tags & re-insert fields │
                         │               └───────────────┬──────────────┘
                         │                               │
                         ▼                               ▼
                  ┌──────────────────────────────────────────────┐
                  │          PASSED TO DETERMINISTIC PARSER      │
                  └──────────────────────────────────────────────┘
\end{verbatim}

\subsubsection{Stage 1: Structural Outline
Prompting}\label{stage-1-structural-outline-prompting}

The generator model is supplied with the citizen's high-level
background, the public administration procedure, and the required
document genre (e.g., formal complaint, petition, judicial declaration).
The model outputs a structured outline specifying the narrative
progression, sections, and disclosure justifications.

\subsubsection{Stage 2: Narrative Prose
Drafting}\label{stage-2-narrative-prose-drafting}

In the second stage, the structural outline is expanded into full
natural language text. The prompt imposes strict operational
constraints: * \textbf{Verbatim Locking}: Atomic identifiers (CCCD,
phone number, address) must appear exactly as specified in the manifest.
* \textbf{Semantic Tagging}: Free-form sensitive disclosures lacking
fixed formats must be wrapped in lightweight semantic delimiters:
\[\texttt{⟦health\_status⟧} \, \text{chẩn đoán suy thận giai đoạn 3} \, \texttt{⟦/health\_status⟧}\]
* \textbf{Distractor Integration}: Document numbers and organizational
titles must be seamlessly woven into official headers.

\subsubsection{Stage 2b: Syntax Revision and Coverage
Validation}\label{stage-2b-syntax-revision-and-coverage-validation}

The generated text \(X_{\text{tagged}}\) is audited by
\texttt{coverage()} (\texttt{core/annotation.py}): 1. \textbf{Manifest
Coverage Check}: Verifies that all mandatory manifest entities appear
within the text. 2. \textbf{Tag Syntax Validation}: Verifies that all
inline delimiters are balanced and conform to the regular expression
\texttt{⟦{[}a-zA-Z\_{]}+⟧...⟦/{[}a-zA-Z\_{]}+⟧}.

If an entity is omitted or a tag is malformed, the pipeline invokes an
automated \textbf{Revision Loop} at low decoding temperature
(\(T = 0.20\)), presenting the model with targeted diagnostic feedback.
The loop executes up to three retries before discarding non-compliant
generations, achieving an overall manifest compliance rate exceeding
\(98.5\%\).

\begin{center}\rule{0.5\linewidth}{0.5pt}\end{center}

\subsection{4.6. Deterministic Anchor-Tag
Parsing}\label{deterministic-anchor-tag-parsing}

\begin{verbatim}
Algorithm 2: Single-Pass Linear-Time LIFO Stack Parser for Inline Markup
──────────────────────────────────────────────────────────────────────────────────
Input  : Tagged text string X_tagged, Sensitivity map SENS
Output : Normalized clean text X_clean, Extracted span set S_anchor

1:  clean_chars ← []
2:  S_anchor    ← ∅
3:  stack       ← []   // LIFO stack storing tuples of (field_name, start_offset)
4:  i ← 0, n ← |X_tagged|
5:  while i < n do
6:      if X_tagged[i] = "⟦" then
7:          close_idx ← FindNext("⟧", X_tagged, from = i + 1)
8:          if close_idx ≠ -1 then
9:              tag_content ← Substring(X_tagged, i + 1, close_idx)
10:             if tag_content begins with "/" then       // Closing delimiter ⟦/field⟧
11:                 field ← Substring(tag_content, 1)
12:                 if stack ≠ [] ∧ Top(stack).field = field then
13:                     (fld, s_offset) ← Pop(stack)
14:                     e_offset ← |clean_chars|
15:                     span ← InstantiateSpan(s_offset, e_offset, fld, SENS[fld], "anchor_tag")
16:                     S_anchor ← S_anchor ∪ {span}
17:                 end if
18:                 i ← close_idx + 1
19:                 continue
20:             else if IsValidIdentifier(tag_content) then // Opening delimiter ⟦field⟧
21:                 Push(stack, (tag_content, |clean_chars|))
22:                 i ← close_idx + 1
23:                 continue
24:             end if
25:         end if
26:     end if
27:     Append(clean_chars, X_tagged[i])
28:     i ← i + 1
29: end while
30: X_clean ← Join(clean_chars)
31: return (X_clean, S_anchor)
──────────────────────────────────────────────────────────────────────────────────
\end{verbatim}

\subsubsection{Algorithmic Guarantees}\label{algorithmic-guarantees}

\begin{enumerate}
\def\labelenumi{\arabic{enumi}.}
\tightlist
\item
  \textbf{Zero Index Drift}: The start coordinate \(s_i\) is recorded at
  the instantaneous length of the output buffer \texttt{clean\_chars}
  when the opening tag is parsed. The end coordinate \(e_i\) is recorded
  when the closing tag is matched. Because tags are skipped rather than
  copied, the resulting offsets map strictly to code points in
  \(X_{\text{clean}}\).
\item
  \textbf{Complete Delimiter Elimination}: Output text
  \(X_{\text{clean}}\) contains no residual tag markers (\texttt{⟦} or
  \texttt{⟧}), ensuring downstream models are trained and evaluated on
  natural prose.
\end{enumerate}

\begin{center}\rule{0.5\linewidth}{0.5pt}\end{center}

\subsection{4.7. Deterministic Character-Span
Annotation}\label{deterministic-character-span-annotation}

Following anchor tag extraction, the pipeline calibrates ground-truth
spans for the remaining structured fields on \(X_{\text{clean}}\) using
three complementary algorithmic matchers:

\subsubsection{\texorpdfstring{Layer 1: Guarded Verbatim Matching
(\texttt{match\_verbatim})}{Layer 1: Guarded Verbatim Matching (match\_verbatim)}}\label{layer-1-guarded-verbatim-matching-match_verbatim}

Applied to atomic identifiers (\texttt{cccd}, \texttt{phone},
\texttt{email}, \texttt{passport\_number}, \texttt{driver\_license},
\texttt{vehicle\_plate}, \texttt{tax\_code}, \texttt{bank\_account},
\texttt{social\_insurance\_no}, \texttt{health\_insurance\_no},
\texttt{address}, \texttt{full\_name}). The algorithm compiles
Unicode-aware zero-width lookaround boundaries around normalized surface
strings:
\[R(v) = \texttt{(?<!\textbackslash{}w)} + \text{re.escape}(\text{NFC}(v)) + \texttt{(?!\textbackslash{}w)}\]
Boundary assertions prevent substring false matches (e.g., matching a
4-digit tax suffix inside a 10-digit number).

\subsubsection{\texorpdfstring{Layer 2: Cue-Context Disambiguation
(\texttt{match\_cue\_context})}{Layer 2: Cue-Context Disambiguation (match\_cue\_context)}}\label{layer-2-cue-context-disambiguation-match_cue_context}

Applied to polysemous and homographic fields (\texttt{dob},
\texttt{gender}, \texttt{marital\_status}, \texttt{nationality},
\texttt{ethnicity}, \texttt{religion}, \texttt{political\_view},
\texttt{location\_data}, \texttt{behavioral\_data},
\texttt{place\_components}). The matcher identifies candidate surface
strings and scans backwards through a sliding window of at most
\(W = 80\) characters, strictly terminated by sentence boundaries
(\texttt{.}, \texttt{;}, \texttt{\textbackslash{}n}, \texttt{!},
\texttt{?}):
\[\text{Valid}(v) \Longleftrightarrow \exists c \in \mathcal{C}_{\text{cue}}(f) \quad \text{in clause segment preceding } v\]
This mechanism reliably differentiates demographic declarations
(\emph{``dân tộc: Kinh''}) from identical common nouns (\emph{``kinh
tế''}), eliminating lexical false positives.

\subsubsection{\texorpdfstring{Layer 3: Disjoint Name Component
Decomposition
(\texttt{split\_vietnamese\_name})}{Layer 3: Disjoint Name Component Decomposition (split\_vietnamese\_name)}}\label{layer-3-disjoint-name-component-decomposition-split_vietnamese_name}

Applied to citizen personal names. Rather than creating overlapping
annotations, the pipeline identifies the constituent tokens of
\texttt{full\_name} and computes disjoint sub-spans for
\texttt{family\_name}, \texttt{middle\_name}, and \texttt{given\_name}.
This preserves spatial partition properties while enabling fine-grained
onomastic evaluation.

\begin{center}\rule{0.5\linewidth}{0.5pt}\end{center}

\subsection{4.8. Priority Lattice Conflict
Resolution}\label{priority-lattice-conflict-resolution}

To reconcile overlapping candidate spans produced across different
extraction mechanisms, the pipeline applies the \textbf{5-Tier Priority
Lattice} (\(\text{PRIORITIES}\)) formalized in Section 3.5.

\begin{verbatim}
       Candidate Spans S_cand
     (Verbatim, Cue, Name, Anchor,
      and doc_code distractors)
                 │
                 ▼
     Sort by: Priority Weight (desc)
              Span Length (desc)
              Start Offset (asc)
                 │
                 ▼
     Greedy Non-Overlapping Selection
                 │
                 ├─► [Tier 5: Interval Subtraction if overlaps higher tier]
                 │
                 ▼
     Resolved Non-Overlapping Spans S_resolved
                 │
                 ▼
     Filter out Distractors (label == 'O')
                 │
                 ▼
     Final Ground Truth Y* in dataset.jsonl
\end{verbatim}

\subsubsection{Tie-Breaking and Interval Trimming
Execution}\label{tie-breaking-and-interval-trimming-execution}

\begin{enumerate}
\def\labelenumi{\arabic{enumi}.}
\tightlist
\item
  \textbf{Candidate Sorting}: All candidate spans
  \(\mathcal{S}_{\text{cand}}\) are sorted by:
  \[\text{key} = \bigl(-\text{weight}(s), -(s.\text{end} - s.\text{start}), s.\text{start}\bigr)\]
\item
  \textbf{Greedy Reservation}: Tiers 1--4 are processed greedily. If a
  candidate span intersects an already accepted span, it is discarded.
\item
  \textbf{Interval Subtraction for Tier 5}: When a Tier 5 narrative
  context (such as \texttt{private\_life}) intersects accepted Tier 1
  atomic credentials (such as a nested \texttt{bank\_account}), the
  pipeline subtracts the occupied interval:
  \[\text{Res}(I) = [s_{\text{broad}}, e_{\text{broad}}) \setminus \bigcup_{k} [s_k, e_k)\]
  Sub-spans satisfying length threshold \(\ge 5\) code points are
  retained as trimmed contextual spans.
\item
  \textbf{Distractor Elimination}: Spans matching adversarial distractor
  codes (\texttt{label:\ "O"}) are utilized during lattice arbitration
  to suppress overlapping erroneous hypotheses, and are then filtered
  out, yielding the finalized benchmark span set \(\mathcal{Y}^*\).
\end{enumerate}

\begin{center}\rule{0.5\linewidth}{0.5pt}\end{center}

\subsection{4.9. Parallel Sanitization Dataset Output
Format}\label{parallel-sanitization-dataset-output-format}

\begin{verbatim}
Figure 4.2: Structured Specimen of an Annotated Parallel Document Instance in ViPII
┌────────────────────────────────────────────────────────────────────────────────────────────────────────┐
│ DOCUMENT METADATA                                                                                      │
│   doc_id: vipii_doc_038192    profile_id: pf_012845    scenario_id: sc_medical_subsidy_04             │
│   form_id: form_mxh_0921      register: administrative_petition                                        │
├────────────────────────────────────────────────────────────────────────────────────────────────────────┤
│ 1. RAW NATURAL LANGUAGE TEXT (X_raw ≡ X_clean)                                                         │
│   "Kính gửi UBND Phường Mai Dịch. Tôi tên là Nguyễn Văn Bình, sinh ngày 12/04/1985, số CCCD          │
│    001085012345, cư trú tại Số 15 ngõ 105 Doãn Kế Thiện. Hiện nay tôi mắc suy thận mạn giai đoạn 3,    │
│    hoàn cảnh gia đình đơn thân nuôi mẹ già 82 tuổi..."                                                 │
├────────────────────────────────────────────────────────────────────────────────────────────────────────┤
│ 2. GROUND-TRUTH CHARACTER SPANS (Y*) [0-indexed Unicode code points]                                   │
│   • [42,  57)  full_name        (PII)  via: verbatim     value: "Nguyễn Văn Bình"                      │
│   • [69,  79)  dob              (PII)  via: cue_context  value: "12/04/1985"                           │
│   • [89, 101)  cccd             (PII)  via: verbatim     value: "001085012345"                         │
│   • [115, 150) address          (PII)  via: verbatim     value: "Số 15 ngõ 105 Doãn Kế Thiện"          │
│   • [165, 194) health_status    (SPI)  via: anchor_tag   value: "mắc suy thận mạn giai đoạn 3"         │
│   • [215, 249) family_relations (SPI)  via: anchor_tag   value: "đơn thân nuôi mẹ già 82 tuổi"         │
├────────────────────────────────────────────────────────────────────────────────────────────────────────┤
│ 3. PARALLEL SANITIZED TEXT: TAG MASKING (X_sanitized, Task 2)                                          │
│   "Kính gửi UBND Phường Mai Dịch. Tôi tên là [HỌ_TÊN], sinh ngày [NGÀY_SINH], số CCCD                 │
│    [CCCD], cư trú tại [ĐỊA_CHỈ]. Hiện nay tôi [TÌNH_TRẠNG_SỨC_KHỎE], hoàn cảnh gia đình               │
│    [QUAN_HỆ_GIA_ĐÌNH]..."                                                                              │
├────────────────────────────────────────────────────────────────────────────────────────────────────────┤
│ 4. PARALLEL SANITIZED TEXT: CONSISTENT SYNTHETIC REPLACEMENT (X_sanitized, Task 2)                    │
│   "Kính gửi UBND Phường Mai Dịch. Tôi tên là Trần Đình Trọng, sinh ngày 28/09/1988, số CCCD           │
│    079088009876, cư trú tại Số 48 đường Cách Mạng Tháng 8. Hiện nay tôi đang điều trị thoái hóa cột    │
│    sống nặng, hoàn cảnh gia đình vợ chồng nuôi 2 con nhỏ..."                                           │
└────────────────────────────────────────────────────────────────────────────────────────────────────────┘
\end{verbatim}

\subsubsection{Downstream Utility of Output
Fields}\label{downstream-utility-of-output-fields}

\begin{itemize}
\tightlist
\item
  \texttt{content}: Represents \(X_{\text{raw}}\)
  (\(X_{\text{clean}}\)), providing raw natural language for Task 1
  (Span-Level Detection).
\item
  \texttt{spans}: Represents ground truth \(\mathcal{Y}^*\), providing
  exact Unicode code-point boundaries for entity evaluation.
\item
  \texttt{sanitized\_tag\_masking}: Ground-truth target for Tag-Masking
  Generative Sanitization (Task 2).
\item
  \texttt{sanitized\_synthetic\_replacement}: Ground-truth target for
  Synthetic Surrogate De-identification (Task 2).
\item
  \texttt{meta}: Audit trail recording demographic origin, procedural
  scenario, and administrative form metadata, enabling
  out-of-distribution split stratification and ablation analysis in
  Sections 5 and 6.
\end{itemize}

By constructing the dataset through this constraint-first, decoupled
architecture, ViPII guarantees zero coordinate drift, comprehensive
statutory coverage of Decree 13, and full reproducibility for Vietnamese
privacy research.

\begin{center}\rule{0.5\linewidth}{0.5pt}\end{center}

\section{Section 5: Dataset Statistics and Quality
Analysis}\label{section-5-dataset-statistics-and-quality-analysis}

In this section, we present a comprehensive empirical evaluation of the
\textbf{ViPII} corpus. We analyze corpus scale and field composition
across the 34 statutory categories, establish lexical and semantic
diversity through multi-scale metric suites (MATTR across varying window
sizes and dense Vendi Scores), and report rigorous verification
diagnostics on annotation integrity, manifest coverage, and span
calibration quality.

\begin{center}\rule{0.5\linewidth}{0.5pt}\end{center}

\subsection{5.1. Corpus Scale and
Composition}\label{corpus-scale-and-composition}

\subsubsection{Global Corpus
Characteristics}\label{global-corpus-characteristics}

The finalized ViPII corpus comprises \textbf{47,881 documents}
containing \textbf{474,874 annotated character spans}. Table 5.1
summarizes the global scale, demographic source assets, register
diversity, and structural length distributions of the benchmark.

\paragraph{Table 5.1: Overall Structural Characteristics of the ViPII
Benchmark
Corpus}\label{table-5.1-overall-structural-characteristics-of-the-vipii-benchmark-corpus}

{\def\LTcaptype{none} % do not increment counter
\begin{longtable}[]{@{}
  >{\raggedright\arraybackslash}p{(\linewidth - 4\tabcolsep) * \real{0.3077}}
  >{\centering\arraybackslash}p{(\linewidth - 4\tabcolsep) * \real{0.3846}}
  >{\raggedright\arraybackslash}p{(\linewidth - 4\tabcolsep) * \real{0.3077}}@{}}
\toprule\noalign{}
\begin{minipage}[b]{\linewidth}\raggedright
Metric / Dimension
\end{minipage} & \begin{minipage}[b]{\linewidth}\centering
Statistical Value
\end{minipage} & \begin{minipage}[b]{\linewidth}\raggedright
Operational Description \& Sub-Categorization
\end{minipage} \\
\midrule\noalign{}
\endhead
\bottomrule\noalign{}
\endlastfoot
\textbf{Total Documents (\(D\))} & \textbf{47,881} & Fully realized,
clean natural language documents (\(X_{\text{raw}}\)). \\
\textbf{Total Character Spans (\(M\))} & \textbf{474,874} & Flat,
non-overlapping ground-truth annotations (\(\mathcal{Y}^*\)). \\
\textbf{Average Spans per Document} & \textbf{9.92} (median: 10) &
Range: 4 to 22 spans per document. \\
\textbf{Statutory Data Distribution} & & \\
• Basic Personal Data (\textbf{PII}) & \textbf{301,545} (63.50\%) &
High-frequency civil registration credentials and contact fields. \\
• Sensitive Personal Data (\textbf{SPI}) & \textbf{173,329} (36.50\%) &
High-risk personal disclosures (health, judicial, financial, family). \\
\textbf{Foundational Asset Coverage} & & \\
• Demographic Profile Bank & \textbf{70,000} profiles & Artificially
generated citizens across 54 ethnic groups. \\
• Situational Scenarios & \textbf{19} scenarios & Procedural motivations
across civil, labor, health, and judicial domains. \\
• Administrative Forms & \textbf{9,020} forms & Reference metadata
mapped across 8 ministerial domains. \\
• Public Service Domains & \textbf{8} domains & Justice (24\%), Public
Security (22\%), Health (18\%), Labor/Social (14\%), Transport (10\%),
Finance (6\%), Education (4\%), Construction (2\%). \\
\textbf{Stylistic Register Breakdown} & & \\
• Formal Administrative Dossiers & \textbf{21,546} (45.00\%) & Official
petitions, administrative applications, and declarations. \\
• First-Person Narrative Petitions & \textbf{14,364} (30.00\%) &
Explanatory statements, hardship appeals, and judicial narratives. \\
• Citizen-Official Dialogues & \textbf{7,182} (15.00\%) & Consultative
public service dialogues and intake interviews. \\
• Official Meeting Minutes \& Records & \textbf{4,789} (10.00\%) & Case
notes, investigative summaries, and verification reports. \\
\textbf{Document Length Statistics} & & \\
• Total Token Count (Syllables) & \textbf{13,406,680} & Monosyllabic
morphemes after NFC normalization. \\
• Mean Document Length & \textbf{280.0} syllables & Median: 274.0
syllables (Std: 48.5; Range: 104 to 586). \\
• Vocabulary Size (\(|\mathcal{V}|\)) & \textbf{38,420} types & Unique
lexical types across the corpus. \\
\textbf{Generator Architectures} & & \\
• Training Generators & 85\% of corpus & DeepSeek-V4-Flash,
Gemini-2.5-Flash, Qwen-2.5-72B. \\
• Evaluation Generator (OOD) & 15\% of corpus & Mistral-Small,
GPT-4o-Mini (Isolated model family split). \\
\end{longtable}
}

\begin{center}\rule{0.5\linewidth}{0.5pt}\end{center}

\subsubsection{Granular Distribution of the 34
Fields}\label{granular-distribution-of-the-34-fields}

Table 5.2 provides the complete frequency census and character span
length statistics across all 34 fields, highlighting the structural
contrast between compact alphanumeric PII and variable-length narrative
SPI.

\paragraph{Table 5.2: Granular Frequency Census and Span Length
Distribution across the 34
Fields}\label{table-5.2-granular-frequency-census-and-span-length-distribution-across-the-34-fields}

{\def\LTcaptype{none} % do not increment counter
\begin{longtable}[]{@{}
  >{\centering\arraybackslash}p{(\linewidth - 16\tabcolsep) * \real{0.1136}}
  >{\raggedright\arraybackslash}p{(\linewidth - 16\tabcolsep) * \real{0.0909}}
  >{\centering\arraybackslash}p{(\linewidth - 16\tabcolsep) * \real{0.1136}}
  >{\centering\arraybackslash}p{(\linewidth - 16\tabcolsep) * \real{0.1136}}
  >{\centering\arraybackslash}p{(\linewidth - 16\tabcolsep) * \real{0.1136}}
  >{\centering\arraybackslash}p{(\linewidth - 16\tabcolsep) * \real{0.1136}}
  >{\centering\arraybackslash}p{(\linewidth - 16\tabcolsep) * \real{0.1136}}
  >{\centering\arraybackslash}p{(\linewidth - 16\tabcolsep) * \real{0.1136}}
  >{\centering\arraybackslash}p{(\linewidth - 16\tabcolsep) * \real{0.1136}}@{}}
\toprule\noalign{}
\begin{minipage}[b]{\linewidth}\centering
STT
\end{minipage} & \begin{minipage}[b]{\linewidth}\raggedright
Field Identifier (\(f\))
\end{minipage} & \begin{minipage}[b]{\linewidth}\centering
Tier
\end{minipage} & \begin{minipage}[b]{\linewidth}\centering
Primary Mechanism
\end{minipage} & \begin{minipage}[b]{\linewidth}\centering
Span Count
\end{minipage} & \begin{minipage}[b]{\linewidth}\centering
Relative \%
\end{minipage} & \begin{minipage}[b]{\linewidth}\centering
Mean Len (chars)
\end{minipage} & \begin{minipage}[b]{\linewidth}\centering
Median Len
\end{minipage} & \begin{minipage}[b]{\linewidth}\centering
Min--Max Len
\end{minipage} \\
\midrule\noalign{}
\endhead
\bottomrule\noalign{}
\endlastfoot
1 & \texttt{full\_name} & PII & Verbatim & 45,210 & 9.52\% & 15.4 & 15 &
7 -- 28 \\
2 & \texttt{cccd} & PII & Verbatim & 38,914 & 8.19\% & 12.0 & 12 & 12 --
14 \\
3 & \texttt{phone} & PII & Verbatim & 36,420 & 7.67\% & 10.2 & 10 & 10
-- 12 \\
4 & \texttt{address} & PII & Verbatim & 35,112 & 7.39\% & 46.8 & 45 & 22
-- 88 \\
5 & \texttt{dob} & PII & Cue-Context & 34,208 & 7.20\% & 10.0 & 10 & 10
-- 10 \\
6 & \texttt{gender} & PII & Cue-Context & 28,450 & 5.99\% & 3.0 & 3 & 3
-- 4 \\
7 & \texttt{email} & PII & Verbatim & 22,140 & 4.66\% & 22.4 & 22 & 14
-- 38 \\
8 & \texttt{tax\_code} & PII & Verbatim & 14,210 & 2.99\% & 10.0 & 10 &
10 -- 13 \\
9 & \texttt{nationality} & PII & Cue-Context & 13,840 & 2.91\% & 8.0 & 8
& 8 -- 9 \\
10 & \texttt{marital\_status} & PII & Cue-Context & 11,200 & 2.36\% &
7.6 & 8 & 3 -- 12 \\
11 & \texttt{driver\_license} & PII & Verbatim & 9,840 & 2.07\% & 12.0 &
12 & 12 -- 12 \\
12 & \texttt{passport\_number} & PII & Verbatim & 8,920 & 1.88\% & 8.0 &
8 & 8 -- 8 \\
13 & \texttt{vehicle\_plate} & PII & Verbatim & 7,650 & 1.61\% & 9.4 & 9
& 8 -- 11 \\
14 & \texttt{cmnd} & PII & Verbatim & 6,420 & 1.35\% & 9.0 & 9 & 9 --
9 \\
15 & \texttt{digital\_account} & PII & Anchor-Tag & 5,110 & 1.08\% &
18.2 & 17 & 8 -- 34 \\
16 & \texttt{family\_name} & PII & Name-Comp. & 4,210 & 0.89\% & 5.8 & 6
& 2 -- 14 \\
17 & \texttt{middle\_name} & PII & Name-Comp. & 3,920 & 0.83\% & 4.2 & 4
& 2 -- 11 \\
18 & \texttt{given\_name} & PII & Name-Comp. & 3,771 & 0.79\% & 4.6 & 4
& 2 -- 8 \\
\textbf{---} & \textbf{Subtotal Basic PII} & \textbf{PII} & --- &
\textbf{301,545} & \textbf{63.50\%} & \textbf{14.2} & \textbf{12} &
\textbf{2 -- 88} \\
19 & \texttt{ethnicity} & SPI & Cue-Context & 24,180 & 5.09\% & 4.8 & 4
& 3 -- 11 \\
20 & \texttt{health\_status} & SPI & Anchor-Tag & 21,450 & 4.52\% & 58.6
& 56 & 18 -- 142 \\
21 & \texttt{religion} & SPI & Cue-Context & 18,920 & 3.98\% & 8.8 & 9 &
3 -- 16 \\
22 & \texttt{private\_life} & SPI & Anchor-Tag & 17,640 & 3.71\% & 72.4
& 69 & 24 -- 186 \\
23 & \texttt{bank\_account} & SPI & Verbatim & 16,810 & 3.54\% & 14.2 &
14 & 10 -- 20 \\
24 & \texttt{social\_insurance\_no} & SPI & Verbatim & 15,200 & 3.20\% &
10.0 & 10 & 10 -- 10 \\
25 & \texttt{health\_insurance\_no} & SPI & Verbatim & 14,110 & 2.97\% &
15.0 & 15 & 10 -- 15 \\
26 & \texttt{political\_view} & SPI & Cue-Context & 11,240 & 2.37\% &
16.4 & 18 & 7 -- 26 \\
27 & \texttt{family\_relations} & SPI & Anchor-Tag & 9,840 & 2.07\% &
44.2 & 42 & 14 -- 96 \\
28 & \texttt{place\_components} & SPI & Cue-Context & 6,820 & 1.44\% &
24.8 & 24 & 12 -- 48 \\
29 & \texttt{criminal\_record} & SPI & Anchor-Tag & 5,420 & 1.14\% &
48.6 & 46 & 18 -- 112 \\
30 & \texttt{location\_data} & SPI & Cue-Context & 3,920 & 0.83\% & 28.4
& 28 & 14 -- 54 \\
31 & \texttt{eid\_credentials} & SPI & Anchor-Tag & 2,840 & 0.60\% &
22.0 & 22 & 16 -- 28 \\
32 & \texttt{biometric} & SPI & Anchor-Tag & 2,110 & 0.44\% & 38.6 & 36
& 16 -- 82 \\
33 & \texttt{behavioral\_data} & SPI & Cue-Context & 1,810 & 0.38\% &
18.2 & 18 & 11 -- 32 \\
34 & \texttt{sexual\_orientation} & SPI & Anchor-Tag & 1,009 & 0.21\% &
42.1 & 39 & 21 -- 88 \\
\textbf{---} & \textbf{Subtotal Sensitive SPI} & \textbf{SPI} & --- &
\textbf{173,329} & \textbf{36.50\%} & \textbf{31.4} & \textbf{22} &
\textbf{3 -- 186} \\
\textbf{ALL} & \textbf{Total Benchmark} & --- & --- & \textbf{474,874} &
\textbf{100.00\%} & \textbf{20.5} & \textbf{12} & \textbf{2 -- 186} \\
\end{longtable}
}

\begin{verbatim}
                                  PII AND SPI FIELD FREQUENCY CENSUS
  full_name [PII]         ████████████████████████ 45,210 (9.52%)
  cccd [PII]              ████████████████████ 38,914 (8.19%)
  phone [PII]             ███████████████████ 36,420 (7.67%)
  address [PII]           ██████████████████ 35,112 (7.39%)
  dob [PII]               █████████████████ 34,208 (7.20%)
  gender [PII]            ██████████████ 28,450 (5.99%)
  ethnicity [SPI]         ████████████ 24,180 (5.09%)
  email [PII]             ███████████ 22,140 (4.66%)
  health_status [SPI]     ███████████ 21,450 (4.52%)
  religion [SPI]          █████████ 18,920 (3.98%)
  private_life [SPI]      █████████ 17,640 (3.71%)
  bank_account [SPI]      ████████ 16,810 (3.54%)
  social_insur [SPI]      ███████ 15,200 (3.20%)
  health_insur [SPI]      ███████ 14,110 (2.97%)
  tax_code [PII]          ███████ 14,210 (2.99%)
  nationality [PII]       ███████ 13,840 (2.91%)
  political_view [SPI]    ██████ 11,240 (2.37%)
  marital_status [PII]    ██████ 11,200 (2.36%)
  family_relations [SPI]  █████ 9,840 (2.07%)
  driver_license [PII]    █████ 9,840 (2.07%)
  passport_number [PII]   ████ 8,920 (1.88%)
  vehicle_plate [PII]     ████ 7,650 (1.61%)
  place_components [SPI]  ███ 6,820 (1.44%)
  cmnd [PII]              ███ 6,420 (1.35%)
  criminal_record [SPI]   ███ 5,420 (1.14%)
  digital_account [PII]   ███ 5,110 (1.08%)
  family_name [PII]       ██ 4,210 (0.89%)
  location_data [SPI]     ██ 3,920 (0.83%)
  middle_name [PII]       ██ 3,920 (0.83%)
  given_name [PII]        ██ 3,771 (0.79%)
  eid_credentials [SPI]   █ 2,840 (0.60%)
  biometric [SPI]         █ 2,110 (0.44%)
  behavioral_data [SPI]   █ 1,810 (0.38%)
  sexual_orientation [SPI]▏ 1,009 (0.21%)
\end{verbatim}

\emph{Figure 5.1: Distribution and frequency census of character spans
across all 34 fields in the ViPII corpus, illustrating the high density
of mandatory civil credentials alongside a substantial tail of
high-liability sensitive disclosures.}

\begin{center}\rule{0.5\linewidth}{0.5pt}\end{center}

\subsection{5.2. Lexical Diversity, Semantic Representation, and
Anti-Mode-Collapse
Analysis}\label{lexical-diversity-semantic-representation-and-anti-mode-collapse-analysis}

A critical concern in synthetic text generation is whether the corpus
exhibits authentic linguistic diversity or collapses into formulaic
templates (\emph{mode collapse}). To address this, we evaluate ViPII
across multi-window lexical turnover (MATTR), semantic feature
dispersion (Vendi Score), document length distributions, and n-gram
near-duplicate analysis.

\subsubsection{Multi-Window Lexical Turnover: MATTR
Evaluation}\label{multi-window-lexical-turnover-mattr-evaluation}

In isolating languages such as Vietnamese, monosyllabic morphemes
(\emph{tiếng}) are whitespace-delimited. Evaluating diversity solely at
an ultra-narrow window (\(w=100\)) introduces an intra-document
artifact, as administrative formalities naturally avoid local
repetition. To provide an uncompromised evaluation, Table 5.3 reports
Moving-Average Type-Token Ratio (\(\text{MATTR}\)) across expanding
window sizes (\(w \in \{10, 20, 50, 100\}\)), accompanied by
corpus-level Measure of Textual Lexical Diversity (\(\text{MTLD}\)) and
vocabulary turnover metrics.

\[\text{MATTR}_w(D) = \frac{1}{N - w + 1} \sum_{i=1}^{N - w + 1} \frac{|\text{UniqueTokens}(D[i : i+w-1])|}{w}\]

\paragraph{Table 5.3: Multi-Window Lexical Diversity and Semantic
Dispersion across ViPII
Subsets}\label{table-5.3-multi-window-lexical-diversity-and-semantic-dispersion-across-vipii-subsets}

{\def\LTcaptype{none} % do not increment counter
\begin{longtable}[]{@{}
  >{\raggedright\arraybackslash}p{(\linewidth - 12\tabcolsep) * \real{0.1176}}
  >{\centering\arraybackslash}p{(\linewidth - 12\tabcolsep) * \real{0.1471}}
  >{\centering\arraybackslash}p{(\linewidth - 12\tabcolsep) * \real{0.1471}}
  >{\centering\arraybackslash}p{(\linewidth - 12\tabcolsep) * \real{0.1471}}
  >{\centering\arraybackslash}p{(\linewidth - 12\tabcolsep) * \real{0.1471}}
  >{\centering\arraybackslash}p{(\linewidth - 12\tabcolsep) * \real{0.1471}}
  >{\centering\arraybackslash}p{(\linewidth - 12\tabcolsep) * \real{0.1471}}@{}}
\toprule\noalign{}
\begin{minipage}[b]{\linewidth}\raggedright
Data Subset / Generator Split
\end{minipage} & \begin{minipage}[b]{\linewidth}\centering
\(\text{MATTR}_{w=10}\)
\end{minipage} & \begin{minipage}[b]{\linewidth}\centering
\(\text{MATTR}_{w=20}\)
\end{minipage} & \begin{minipage}[b]{\linewidth}\centering
\(\text{MATTR}_{w=50}\)
\end{minipage} & \begin{minipage}[b]{\linewidth}\centering
\(\text{MATTR}_{w=100}\)
\end{minipage} & \begin{minipage}[b]{\linewidth}\centering
\(\text{MTLD}\) Factor
\end{minipage} & \begin{minipage}[b]{\linewidth}\centering
Dense Vendi Score (\(V_{\text{dense}}\))
\end{minipage} \\
\midrule\noalign{}
\endhead
\bottomrule\noalign{}
\endlastfoot
\textbf{Formal Administrative Dossiers} & 0.982 & 0.941 & 0.868 & 0.794
& 74.2 & 84.5 \\
\textbf{Narrative Petitions \& Appeals} & 0.988 & 0.954 & 0.892 & 0.835
& 92.6 & 118.2 \\
\textbf{Citizen-Official Dialogues} & 0.991 & 0.962 & 0.910 & 0.852 &
104.1 & 132.4 \\
\textbf{Official Meeting Records} & 0.979 & 0.935 & 0.854 & 0.781 & 68.4
& 76.8 \\
\textbf{Generator: DeepSeek-V4-Flash} & 0.985 & 0.948 & 0.880 & 0.821 &
86.4 & 105.4 \\
\textbf{Generator: Gemini-2.5-Flash} & 0.986 & 0.951 & 0.886 & 0.828 &
89.2 & 110.6 \\
\textbf{Generator: Qwen-2.5-72B} & 0.984 & 0.946 & 0.878 & 0.819 & 84.8
& 102.1 \\
\textbf{Held-Out OOD (Mistral-Small)} & 0.989 & 0.956 & 0.895 & 0.839 &
94.5 & 124.0 \\
\textbf{Overall ViPII Corpus} & \textbf{0.985} & \textbf{0.948} &
\textbf{0.881} & \textbf{0.822} & \textbf{88.1} & \textbf{108.7} \\
\end{longtable}
}

\paragraph{Analysis of Metric
Trajectory}\label{analysis-of-metric-trajectory}

As the window expands from \(w=10\) to \(w=100\), MATTR scales
gracefully from \(0.985\) to \(0.822\). Even under \(w=100\) (covering
approximately 50--60 words, equivalent to 2--3 full clauses), lexical
turnover remains above \(0.82\), confirming that administrative texts
maintain continuous vocabulary variation rather than cycling repetitive
stock phrases.

\subsubsection{Dense Semantic Dispersion: Vendi Score
Analysis}\label{dense-semantic-dispersion-vendi-score-analysis}

Rather than relying on sparse TF-IDF vectors (which inflate cosine
orthogonality due to high-entropy unique numbers like CCCD and phone
numbers), we calculate the \textbf{Vendi Score using 1,024-dimensional
dense multilingual embeddings} (\(\text{text-embedding-3-large}\)):
\[V = \exp\left(-\sum_{i=1}^n \lambda_i \ln \lambda_i\right)\] where
\(\lambda_i\) are the eigenvalues of the normalized kernel matrix
\(K/n\). The overall corpus achieves a dense Vendi Score of
\textbf{\(V_{\text{dense}} = 108.7\)}. Because the situational scenario
catalog contains 19 canonical anchors, a dense Vendi Score of \(108.7\)
demonstrates that semantic variation within each scenario spans on
average \(5.7\) distinct semantic sub-clusters, confirming that
documents do not collapse into rigid templates.

\begin{verbatim}
                  DOCUMENT LENGTH FREQUENCY HISTOGRAM
  [100 - 150]   ██ 1,240 (2.59%)
  [151 - 200]   ██████ 4,820 (10.07%)
  [201 - 250]   ██████████████████ 14,650 (30.59%)
  [251 - 300]   █████████████████████ 17,120 (35.75%)
  [301 - 350]   █████████ 7,450 (15.56%)
  [351 - 400]   ██ 1,840 (3.84%)
  [401 - 450]   █ 560 (1.17%)
  [451 - 600]   ▏ 201 (0.42%)
\end{verbatim}

\emph{Figure 5.2: Document length distribution (syllables) across the
ViPII corpus, displaying a Gaussian profile centered at 280 syllables,
matching the length of authentic Vietnamese administrative dossiers.}

\subsubsection{Near-Duplicate and Template Surface Overlap
Analysis}\label{near-duplicate-and-template-surface-overlap-analysis}

To verify that documents do not mirror identical phrasing: 1.
\textbf{MinHash LSH De-duplication}: Using 128 permutation hashes over
character 5-grams, pairwise similarity analysis revealed that
\textbf{less than 1.14\%} of document pairs exhibited Jaccard similarity
exceeding \(0.60\). 2. \textbf{Structural Contrast (Atomic PII
vs.~Narrative SPI)}: As established in Table 5.2, atomic civil
credentials exhibit invariant lengths (CCCD: 12 chars; phone: 10 chars;
tax code: 10 chars), whereas sensitive narrative clauses span broad
syntactic distributions (\texttt{health\_status}: mean 58.6 chars, max
142 chars; \texttt{private\_life}: mean 72.4 chars, max 186 chars). This
contrast forces downstream neural models to balance precision on exact
numeric patterns with contextual semantic boundaries on natural language
prose.

\begin{center}\rule{0.5\linewidth}{0.5pt}\end{center}

\subsection{5.3. Annotation Integrity and Validation
Verification}\label{annotation-integrity-and-validation-verification}

To ensure absolute reproducibility and scientific validity, every
instance in ViPII is audited through automated algorithmic verification
pipelines accompanied by expert human spot-checking.

\paragraph{Table 5.4: Annotation Integrity, Algorithmic Verification,
and Quality
Diagnostics}\label{table-5.4-annotation-integrity-algorithmic-verification-and-quality-diagnostics}

{\def\LTcaptype{none} % do not increment counter
\begin{longtable}[]{@{}
  >{\raggedright\arraybackslash}p{(\linewidth - 6\tabcolsep) * \real{0.2222}}
  >{\centering\arraybackslash}p{(\linewidth - 6\tabcolsep) * \real{0.2778}}
  >{\centering\arraybackslash}p{(\linewidth - 6\tabcolsep) * \real{0.2778}}
  >{\raggedright\arraybackslash}p{(\linewidth - 6\tabcolsep) * \real{0.2222}}@{}}
\toprule\noalign{}
\begin{minipage}[b]{\linewidth}\raggedright
Quality Metric / Verification Dimension
\end{minipage} & \begin{minipage}[b]{\linewidth}\centering
Target Standard
\end{minipage} & \begin{minipage}[b]{\linewidth}\centering
Measured Value
\end{minipage} & \begin{minipage}[b]{\linewidth}\raggedright
Diagnostic Outcome \& Validation Policy
\end{minipage} \\
\midrule\noalign{}
\endhead
\bottomrule\noalign{}
\endlastfoot
\textbf{Manifest Coverage Rate} & \(\ge 98.0\%\) & \textbf{98.62\%} &
Proportion of mandatory manifest entities present in clean text. \\
\textbf{Tag Syntax Validity (Initial Draft)} & \(\ge 95.0\%\) &
\textbf{96.40\%} & Proportion of generations with balanced, valid
\texttt{⟦field⟧...⟦/field⟧} markup. \\
\textbf{Tag Syntax Validity (Post-Revision)} & \(\ge 99.0\%\) &
\textbf{99.42\%} & Yield after executing the automated Stage 2b revision
loop (\(T=0.20\)). \\
\textbf{Substring Exact-Match Integrity} & \(100.0\%\) &
\textbf{100.00\%} & Zero-tolerance invariant:
\(X_{\text{clean}}[s_i:e_i] \equiv \text{surface}(y_i^*)\) across all
spans. \\
\textbf{Initial Pipeline Validation Pass Rate} & \(\ge 95.0\%\) &
\textbf{96.78\%} & Accepted on initial inference pass; 3.22\% routed to
revision/regen. \\
\textbf{Candidate Overlap Collision Rate} & Baseline & \textbf{14.82\%}
& Proportion of candidate spans in \(\mathcal{S}_{\text{cand}}\)
requiring resolution by \(\Omega\). \\
\textbf{Overlap Resolution Success Rate} & \(100.0\%\) &
\textbf{100.00\%} & Algorithm 1 yields zero intersecting spans
(\(e_i^* \le s_j^*\)). \\
\textbf{Adversarial Distractor Rejection Rate} & \(100.0\%\) &
\textbf{100.00\%} & All 12-digit \texttt{doc\_code} tokens (label
\texttt{O}) successfully purged from \(\mathcal{Y}^*\). \\
\textbf{Unicode Code-Point Coordinate Drift} & \(0.0\%\) &
\textbf{0.00\%} & Zero offset deviation between Python string index and
physical text. \\
\textbf{Human Gold Spot-Check Agreement} & \(\ge 95.0\%\) &
\textbf{99.20\%} & Cohen's \(\kappa = 0.941\) across 500 randomly
audited documents. \\
\end{longtable}
}

\subsubsection{Technical Verification
Procedures}\label{technical-verification-procedures}

\begin{enumerate}
\def\labelenumi{\arabic{enumi}.}
\tightlist
\item
  \textbf{Zero-Tolerance Substring Assertion}: Every exported span tuple
  \((s_i, e_i, f_i, l_i)\) is programmatically validated against the raw
  document string:
  \[\text{Assert } X_{\text{clean}}[s_i : e_i] == y_i^*.\text{value} \quad \forall y_i^* \in \mathcal{Y}^*\]
  Over all 474,874 ground-truth spans, zero substring mismatch
  exceptions were encountered.
\item
  \textbf{Candidate Overlap Resolution Efficiency}: Before resolution,
  \(14.82\%\) of candidate spans exhibit boundary collisions
  (predominantly nested family name components inside full names and
  atomic credentials inside broad poverty or health narratives).
  Algorithm 1 deterministically resolves \(100\%\) of collisions,
  converting candidate sets into strict mathematical partitions.
\item
  \textbf{Human Spot-Check Audit}: Two independent native Vietnamese
  annotators audited a stratified sample of \textbf{500 documents}
  (representing 5,124 spans). Annotators evaluated whether extracted
  boundaries accurately captured the complete statutory scope under
  Decree 13. The audit achieved an inter-annotator agreement of
  \textbf{\(\kappa = 0.941\)}, with boundary discrepancies occurring
  exclusively on ambiguous transitional connectives in broad
  \texttt{private\_life} clauses.
\end{enumerate}

The empirical statistics establish the corpus-level basis for the
downstream model evaluation and results reported in Section 6.

\begin{center}\rule{0.5\linewidth}{0.5pt}\end{center}

\section{Section 6: Experimental Evaluation and
Results}\label{section-6-experimental-evaluation-and-results}

This section evaluates compact fine-tuned language models and
general-purpose baselines on the ViPII utility-preserving sanitization
task. We first summarize only the experimental information required to
interpret the comparison, and then focus on the main empirical findings.

\begin{center}\rule{0.5\linewidth}{0.5pt}\end{center}

\subsection{6.1. Concise Experimental
Setup}\label{concise-experimental-setup}

\subsubsection{Data and evaluation
protocol}\label{data-and-evaluation-protocol}

ViPII contains 47,881 documents and 474,874 annotated character spans.
The corpus is partitioned into 37,548 training documents, 4,788
validation documents, and 5,545 test documents. The split is designed to
isolate profiles, administrative forms/templates, scenarios, and
generator families across partitions, reducing identity, template,
contextual, and stylistic leakage. All reported systems are evaluated on
the same held-out test set.

\subsubsection{Compared systems}\label{compared-systems}

We compare eight configurations in four functional groups:

\begin{enumerate}
\def\labelenumi{\arabic{enumi}.}
\tightlist
\item
  \textbf{Proposed compact models:} \texttt{qwen3-0.6b\ (FT)} and
  \texttt{qwen3-1.7b\ (FT)}, supervised fine-tuned on the ViPII training
  partition.
\item
  \textbf{Unadapted compact baselines:} the corresponding
  \texttt{qwen3-0.6b\ (base)} and \texttt{qwen3-1.7b\ (base)} models
  evaluated without ViPII fine-tuning.
\item
  \textbf{General-purpose and reference systems:}
  \texttt{DeepSeek-V4-Flash}, \texttt{Gemini-3.6-Flash-High}, and
  \texttt{GLM-5.1-Flash\ (Pipeline)}. GLM is treated as a reference
  teacher pipeline rather than a directly comparable compact deployment
  model.
\item
  \textbf{Control condition:} untouched raw text, representing maximum
  content retention without privacy sanitization.
\end{enumerate}

All systems receive unannotated Vietnamese text and produce sanitized
text under a common output convention. Training hyperparameters, model
versions, prompts, decoding settings, hardware, and checkpoint-selection
details will be provided in the reproducibility appendix and released
experiment configuration.

\subsubsection{Evaluation metrics}\label{evaluation-metrics}

The primary privacy metric is \textbf{SanRec}, the proportion of
annotated PII/SPI units successfully sanitized. \textbf{SanAtt}
macro-averages sanitization performance across the 34 fields, whereas
\textbf{SanA/R} averages sanitization success across documents. Their
utility counterparts---\textbf{RetRec}, \textbf{RetAtt}, and
\textbf{RetA/R}---measure preservation of non-sensitive operational
content. \textbf{FULL} is the strict document-level success rate: a
document succeeds only when all sensitive units are sanitized and all
evaluated utility units are retained. The formal definitions follow
Section 3.2; implementation-level matching and aggregation rules are
included with the evaluator release.

\begin{center}\rule{0.5\linewidth}{0.5pt}\end{center}

\subsection{6.2. Main Results}\label{main-results}

Table 6.1 reports the central benchmark results. Higher values are
better for every metric.

\subsubsection{Table 6.1: Utility-preserving sanitization results on the
held-out ViPII test set
(\%)}\label{table-6.1-utility-preserving-sanitization-results-on-the-held-out-vipii-test-set}

{\def\LTcaptype{none} % do not increment counter
\begin{longtable}[]{@{}
  >{\raggedright\arraybackslash}p{(\linewidth - 16\tabcolsep) * \real{0.1111}}
  >{\raggedright\arraybackslash}p{(\linewidth - 16\tabcolsep) * \real{0.1111}}
  >{\raggedleft\arraybackslash}p{(\linewidth - 16\tabcolsep) * \real{0.1111}}
  >{\raggedleft\arraybackslash}p{(\linewidth - 16\tabcolsep) * \real{0.1111}}
  >{\raggedleft\arraybackslash}p{(\linewidth - 16\tabcolsep) * \real{0.1111}}
  >{\raggedleft\arraybackslash}p{(\linewidth - 16\tabcolsep) * \real{0.1111}}
  >{\raggedleft\arraybackslash}p{(\linewidth - 16\tabcolsep) * \real{0.1111}}
  >{\raggedleft\arraybackslash}p{(\linewidth - 16\tabcolsep) * \real{0.1111}}
  >{\raggedleft\arraybackslash}p{(\linewidth - 16\tabcolsep) * \real{0.1111}}@{}}
\toprule\noalign{}
\begin{minipage}[b]{\linewidth}\raggedright
System
\end{minipage} & \begin{minipage}[b]{\linewidth}\raggedright
Role
\end{minipage} & \begin{minipage}[b]{\linewidth}\raggedleft
SanAtt
\end{minipage} & \begin{minipage}[b]{\linewidth}\raggedleft
SanA/R
\end{minipage} & \begin{minipage}[b]{\linewidth}\raggedleft
SanRec
\end{minipage} & \begin{minipage}[b]{\linewidth}\raggedleft
RetAtt
\end{minipage} & \begin{minipage}[b]{\linewidth}\raggedleft
RetA/R
\end{minipage} & \begin{minipage}[b]{\linewidth}\raggedleft
RetRec
\end{minipage} & \begin{minipage}[b]{\linewidth}\raggedleft
FULL
\end{minipage} \\
\midrule\noalign{}
\endhead
\bottomrule\noalign{}
\endlastfoot
\textbf{GLM-5.1-Flash (Pipeline)} & Reference teacher & \textbf{96.15} &
\textbf{96.04} & \textbf{85.03} & 98.63 & 98.80 & 97.02 &
\textbf{82.64} \\
\textbf{qwen3-0.6b (FT)} & Proposed compact model & 92.17 & 92.10 &
72.46 & 98.26 & 98.38 & 96.20 & 70.10 \\
\textbf{qwen3-1.7b (FT)} & Proposed compact model & 90.61 & 90.84 &
68.88 & 99.23 & 99.34 & 98.31 & 67.90 \\
\textbf{Gemini-3.6-Flash-High} & Commercial baseline & 89.30 & 90.10 &
68.78 & 97.49 & 97.83 & 94.62 & 64.82 \\
\textbf{DeepSeek-V4-Flash} & General-purpose baseline & 73.72 & 74.37 &
40.97 & 95.84 & 95.96 & 91.90 & 37.72 \\
\textbf{qwen3-0.6b (base)} & Unadapted compact baseline & 19.44 & 19.80
& 4.49 & 80.17 & 81.79 & 70.59 & 0.82 \\
\textbf{qwen3-1.7b (base)} & Unadapted compact baseline & 4.21 & 4.33 &
0.00 & \textbf{99.86} & \textbf{99.84} & \textbf{99.72} & 0.00 \\
\textbf{Untouched raw text} & Control & 1.16 & 1.22 & 0.00 & 99.19 &
99.28 & 98.26 & 0.00 \\
\end{longtable}
}

The reference GLM pipeline obtains the highest overall result, with
85.03\% SanRec and 82.64\% FULL. This result provides an empirical
reference ceiling for the current benchmark. Among the compact
deployment candidates, the fine-tuned 0.6B model achieves the strongest
privacy--utility balance, reaching 72.46\% SanRec and 70.10\% FULL while
retaining 98.26\% of evaluated utility attributes.

\begin{center}\rule{0.5\linewidth}{0.5pt}\end{center}

\subsection{6.3. Effect of ViPII
Fine-Tuning}\label{effect-of-vipii-fine-tuning}

Fine-tuning produces the largest observed performance change. For
Qwen3-0.6B, SanRec increases from 4.49\% to 72.46\% (\textbf{+67.97
percentage points}) and FULL increases from 0.82\% to 70.10\%
(\textbf{+69.28 points}). SanAtt similarly rises from 19.44\% to 92.17\%
(\textbf{+72.73 points}).

The Qwen3-1.7B base model retains nearly all non-sensitive content but
sanitizes none of the evaluated sensitive units, yielding 0\% FULL.
After fine-tuning, it reaches 68.88\% SanRec and 67.90\% FULL while
maintaining 98.31\% RetRec. These paired base-versus-fine-tuned
comparisons indicate that task-specific adaptation, rather than
parameter count alone, is the principal source of sanitization
capability in the evaluated compact models.

\begin{center}\rule{0.5\linewidth}{0.5pt}\end{center}

\subsection{6.4. Compact Models versus General-Purpose
Baselines}\label{compact-models-versus-general-purpose-baselines}

The fine-tuned Qwen3-0.6B model exceeds DeepSeek-V4-Flash by
\textbf{31.49 points in SanRec} (72.46\% versus 40.97\%) and
\textbf{32.38 points in FULL} (70.10\% versus 37.72\%). It also exceeds
Gemini-3.6-Flash-High by \textbf{5.28 points in FULL} (70.10\% versus
64.82\%) and by \textbf{3.68 points in SanRec} (72.46\% versus 68.78\%).

These results show that, under the current ViPII protocol, a compact
model adapted to Vietnamese PII/SPI sanitization can outperform larger
general-purpose baselines on strict end-to-end success. The claim is
limited to the evaluated models, prompts, test set, and metric
definitions; it does not imply universal superiority over frontier
language models.

\begin{center}\rule{0.5\linewidth}{0.5pt}\end{center}

\subsection{6.5. Privacy--Utility Trade-off between the 0.6B and 1.7B
Models}\label{privacyutility-trade-off-between-the-0.6b-and-1.7b-models}

The two fine-tuned models exhibit a consistent trade-off. Qwen3-0.6B
provides stronger sanitization, outperforming Qwen3-1.7B by \textbf{3.58
points in SanRec} and \textbf{2.20 points in FULL}. Conversely,
Qwen3-1.7B preserves more operational content, improving RetAtt from
98.26\% to 99.23\% and RetRec from 96.20\% to 98.31\%.

Accordingly, the 0.6B configuration is the stronger choice when residual
privacy leakage is the primary concern, whereas the 1.7B configuration
may be preferable when conservative content retention is prioritized.
Both models remain below the reference teacher pipeline, leaving
substantial room for improving privacy recall without sacrificing
utility.

\begin{center}\rule{0.5\linewidth}{0.5pt}\end{center}

\subsection{6.6. Summary of Findings}\label{summary-of-findings}

The experimental results support three conclusions. First, unadapted
compact models do not reliably sanitize Vietnamese PII/SPI. Second,
ViPII fine-tuning yields large gains in both privacy recall and strict
document-level success. Third, the 0.6B fine-tuned model offers the best
result among the evaluated compact systems and general-purpose
baselines, while the 1.7B model provides slightly stronger utility
retention. Component ablations, robustness tests, error categories, and
deployment efficiency are examined in the following section.

\end{document}
